\documentclass[11pt,oneside]{article}

\usepackage[T1]{fontenc}
\usepackage{lmodern}
\usepackage{microtype}
\usepackage[margin=1in,headheight=14pt]{geometry}
\usepackage{amsmath,amssymb,amsthm,mathtools}
\usepackage{booktabs,longtable,tabularx,array}
\usepackage{enumitem}
\usepackage{xcolor}
\usepackage{xurl}
\usepackage{hyperref}
\usepackage{listings}
\usepackage{fancyhdr}
\usepackage{titlesec}
\usepackage{needspace}

\definecolor{proofgreen}{RGB}{14,100,65}
\definecolor{openred}{RGB}{155,48,48}
\definecolor{exactviolet}{RGB}{97,56,143}
\definecolor{contextgray}{RGB}{83,92,100}
\definecolor{linkblue}{RGB}{26,91,145}
\definecolor{codegray}{RGB}{247,247,247}

\hypersetup{
  colorlinks=true,
  linkcolor=linkblue,
  citecolor=proofgreen,
  urlcolor=linkblue,
  pdftitle={Alaoglu--Erdos Frontier Update VII: Transported Newton Faces},
  pdfauthor={Research progress report},
  pdfsubject={Unimodular valuation transport, fixed cross-ratios, tangent normal forms, and the remaining cross-prime barrier}
}
\urlstyle{same}

\pagestyle{fancy}
\fancyhf{}
\lhead{Alaoglu--Erd\H{o}s frontier update VII}
\rhead{Transported Newton faces}
\cfoot{\thepage}
\renewcommand{\headrulewidth}{0.4pt}

\titleformat{\section}{\Large\bfseries}{\thesection}{0.75em}{}
\titleformat{\subsection}{\large\bfseries}{\thesubsection}{0.75em}{}
\titleformat{\subsubsection}{\normalsize\bfseries}{\thesubsubsection}{0.75em}{}

\setlist[itemize]{leftmargin=2em,itemsep=0.28em,topsep=0.4em}
\setlist[enumerate]{leftmargin=2.2em,itemsep=0.28em,topsep=0.4em}
\setlength{\parindent}{0pt}
\setlength{\parskip}{0.55em}
\setlength{\emergencystretch}{3.5em}

\newtheorem{theorem}{Theorem}[section]
\newtheorem{proposition}[theorem]{Proposition}
\newtheorem{lemma}[theorem]{Lemma}
\newtheorem{corollary}[theorem]{Corollary}
\newtheorem{conjecture}[theorem]{Conjecture}
\newtheorem{question}[theorem]{Question}
\newtheorem{target}[theorem]{Research target}
\theoremstyle{definition}
\newtheorem{definition}[theorem]{Definition}
\newtheorem{experiment}[theorem]{Exact calibration}
\newtheorem{warning}[theorem]{Caution}
\theoremstyle{remark}
\newtheorem{remark}[theorem]{Remark}

\newcommand{\N}{\mathbb N}
\newcommand{\Nzero}{\mathbb N_0}
\newcommand{\Z}{\mathbb Z}
\newcommand{\Q}{\mathbb Q}
\newcommand{\R}{\mathbb R}
\newcommand{\C}{\mathbb C}
\newcommand{\Abar}{\overline{\mathbb Q}}
\newcommand{\thetaq}{\vartheta}
\newcommand{\eps}{\varepsilon}
\newcommand{\vp}{v_p}
\newcommand{\vell}{v_{\ell}}
\newcommand{\supp}{\operatorname{supp}}
\newcommand{\rank}{\operatorname{rank}}
\newcommand{\Span}{\operatorname{span}}
\newcommand{\conv}{\operatorname{conv}}
\newcommand{\den}{\operatorname{den}}
\newcommand{\num}{\operatorname{num}}
\newcommand{\height}{\operatorname{h}}
\newcommand{\ord}{\operatorname{ord}}
\newcommand{\sgn}{\operatorname{sgn}}
\newcommand{\GL}{\operatorname{GL}}
\newcommand{\PGL}{\operatorname{PGL}}
\newcommand{\SNF}{\operatorname{SNF}}
\newcommand{\AEname}{Alaoglu--Erd\H{o}s}
\newcommand{\provedhere}{\textcolor{proofgreen}{\textbf{[proved in this report]}}}
\newcommand{\conditional}{\textcolor{contextgray}{\textbf{[conditional on a counterexample]}}}
\newcommand{\exactcheck}{\textcolor{exactviolet}{\textbf{[exact symbolic check]}}}
\newcommand{\openstatus}{\textcolor{openred}{\textbf{[open target]}}}
\newcommand{\knownstatus}{\textcolor{contextgray}{\textbf{[classical background]}}}

\newenvironment{statusbox}[1]
{\begin{center}\begin{minipage}{0.94\textwidth}\raggedright\hrule\vspace{0.45em}\textbf{Status:} #1\par\vspace{0.35em}}
{\vspace{0.35em}\hrule\end{minipage}\end{center}}

\lstdefinestyle{pythonstyle}{
  language=Python,
  basicstyle=\ttfamily\scriptsize,
  backgroundcolor=\color{codegray},
  frame=single,
  numbers=left,
  numberstyle=\tiny,
  numbersep=7pt,
  showstringspaces=false,
  breaklines=true,
  columns=fullflexible,
  keepspaces=true,
  tabsize=4
}

\begin{document}
\pagenumbering{gobble}

\begin{titlepage}
\centering
\vspace*{0.7cm}
{\Huge\bfseries The Alaoglu--Erd\H{o}s Conjecture\par}
\vspace{0.25cm}
{\LARGE\bfseries Frontier Update VII\par}
\vspace{0.55cm}
{\LARGE Transported Newton Faces, Fixed Cross-Ratios,\par}
\vspace{0.12cm}
{\LARGE and the Structural-Prime Normal Form\par}
\vspace{0.75cm}
{\large A standalone continuation of the unified report, version 6\par}
\vfill
\begin{minipage}{0.91\textwidth}
\small
\textbf{Bottom line.}
This report does not claim a proof of the conjecture.  It obtains a new exact description of
one of the principal surviving architectures in the supplied unified report: bivariate
near-unit masks and their shared structural-prime Newton faces.  The valuation-depth pair at
any support prime is a unimodular image of that prime's fixed output-valuation vector.  Hence
its consecutive gcd is not merely bounded but exactly constant, every tied Newton edge has a
unique primitive direction of length linear in the convergent denominator, and the monomial
ratio along that edge is a fixed positive rational number independent of the convergent.

This converts local face cancellation into an elementary one-variable problem.  A face
polynomial either has bounded extra valuation in terms of its coefficient height and edge
length, or it contains an explicit candidate-specific tangent binomial.  The latter can be
divided out by a canonical integral normal form without changing the relevant local
valuation.  Thus cancellation at a single structural prime is algebraically removable; any
remaining gain is coefficient-driven or genuinely simultaneous across different prime
slopes.

For two nonproportional structural primes the primitive tangent vectors have length
$\asymp q$ but a determinant independent of $q$.  The resulting tangent lattice has constant
index and extreme condition number.  Exact short-vector identities show why simultaneous
reduction is critical rather than immediately contradictory: a bounded exponent displacement
requires coefficients of size $\asymp q$ in the long tangent basis, reproducing the original
height scale.  This isolates a sharper frontier problem---a cross-prime normal-form
contraction---and removes the entire sublinear-degree and one-prime geometric branch from the
mask search.
\end{minipage}
\vfill
{\large August 24, 2026\par}
\end{titlepage}

\pagenumbering{roman}

\begin{abstract}
Put $\thetaq=\log_2 3$.  The \AEname{} conjecture asserts that $2^x,3^x\in\Z$ implies
$x\in\Z$.  Assume provisionally that a nonintegral solution $\beta$ exists and write
$M=2^\beta$, $A=3^\beta$.  For consecutive convergents $p_k/q_k$ of $\thetaq$, the rational
near-units
\[
 z_k=\frac{A^{q_k}}{M^{p_k}}
\]
converge to $1$.  If $m_\ell=v_\ell(M)$ and $n_\ell=v_\ell(A)$, their local depths are
$u_{\ell,k}=q_kn_\ell-p_km_\ell$.

The first result of this report is the exact transport identity
\[
 (u_{\ell,k},u_{\ell,k+1})
 =(m_\ell,n_\ell)
 \begin{pmatrix}-p_k&-p_{k+1}\\q_k&q_{k+1}\end{pmatrix},
\]
whose matrix is unimodular.  Consequently
$\gcd(u_{\ell,k},u_{\ell,k+1})=\gcd(m_\ell,n_\ell)$, all determinantal divisors of the
multi-prime depth matrix are invariant, and cross-prime determinants are constant.  At a
stable denominator prime define $s_{\ell,k}=-u_{\ell,k}>0$,
$h_\ell=\gcd(m_\ell,n_\ell)$,
$a_{\ell,k}=s_{\ell,k+1}/h_\ell$, and
$b_{\ell,k}=s_{\ell,k}/h_\ell$.  Then
\[
 \frac{z_k^{a_{\ell,k}}}{z_{k+1}^{b_{\ell,k}}}
 =\left(\frac{A^{m_\ell/h_\ell}}{M^{n_\ell/h_\ell}}\right)^{\eps_k},
 \qquad \eps_k=p_{k+1}q_k-p_kq_{k+1}\in\{\pm1\}.
\]
Thus the primitive direction of a tied $\ell$-adic Newton face evaluates at a fixed rational
cross-ratio.  This gives a linear degree threshold for every tied face, an exact
edge-polynomial reduction, a rational-height bound for nonzero edge values, and a canonical
normal form modulo the tangent binomial.  For two prime slopes, the primitive tangent
vectors have a fixed determinant equal to the normalized determinant of the corresponding
output-valuation vectors.  The associated quotient has constant lattice index, but its basis
condition number is $\asymp q_k$; exact resultant identities recover $z_k$ and $z_{k+1}$ as
large powers of the fixed cross-ratios.  This explains both the new simplification and the
remaining obstruction.

All new mathematical statements are proved self-containedly.  A short exact-arithmetic
script checks the sign conventions, gcd invariance, determinant identities, cross-ratios, and
normal-form evaluation on synthetic valuation data.  Announcements by competing systems are
not treated as evidence; the status ledger distinguishes proved results, exact checks,
classical inputs, and open targets.
\end{abstract}

\tableofcontents
\clearpage
\pagenumbering{arabic}

\section{Scope, status, and the precise increment over the unified report}
\label{sec:scope}

\subsection{The problem}

\begin{conjecture}[\AEname{} integer-exponent conjecture]
\label{conj:AE}
For every real $x$,
\[
 2^x\in\Z\quad\text{and}\quad 3^x\in\Z
 \qquad\Longrightarrow\qquad x\in\Z.
\]
\end{conjecture}

The conjecture occurs in the work of Alaoglu and Erd\H{o}s on highly composite numbers
\cite{AE1944}.  It is a specialized $2\times2$ four-exponentials problem: if $x$ is
irrational, the two rows $1,x$ and the two columns $\log2,\log3$ are each linearly independent
over $\Q$, while all four exponentials
\[
 e^{\log2}=2,\qquad e^{\log3}=3,\qquad
 e^{x\log2}=2^x,\qquad e^{x\log3}=3^x
\]
would be algebraic.  Thus the four exponentials conjecture would settle
Conjecture~\ref{conj:AE}; present transcendence theory does not.

The supplied unified report, version 6, already develops a very large corpus: exact
conditional monoid structure, finite verification, determinant and interpolation methods,
local and adelic diagnostics, Kummer and radical reformulations, compact-orbit dynamics,
integer-valued auxiliary functions, Matrix Coefficient specializations, and extensive
negative audits.  Its terminal mask target asks whether shared Newton faces at the fixed
support primes can compress the reduced denominator of a high-contact bivariate polynomial.
The present report addresses exactly that target.  It does not repeat the broad corpus.

\subsection{What is new here}

The increment consists of five tightly connected results.

\begin{enumerate}
\item The consecutive local-depth vector is a \emph{unimodular transform} of the fixed
valuation vector $(v_\ell(M),v_\ell(A))$.  This gives exact gcd and Smith-form invariance.
\item Every tied local Newton face has a primitive tangent vector whose coordinates grow
linearly with the convergent denominators, and whose monomial ratio at $(z_k,z_{k+1})$ is a
fixed rational cross-ratio.
\item A face initial form becomes a one-variable integer polynomial evaluated at that fixed
rational.  Nonzero evaluation has an explicit valuation bound; zero evaluation is exactly a
tangent-binomial factor.
\item There is a canonical integral normal form modulo that tangent binomial.  It removes all
one-prime geometric ties at a coefficient cost proportional to degree divided by the
convergent scale.
\item At two nonproportional primes the tangent vectors form a constant-index, badly
conditioned lattice.  Exact short-vector identities identify the remaining height tax and
show that a successful argument must be cross-prime and coefficient-sensitive.
\end{enumerate}

\begin{statusbox}{All theorem and proposition proofs labeled \provedhere{} are ordinary
paper proofs given in full below.  The report makes no claim that these new statements have
already been translated into or accepted by Lean.  The computational table and script are
\exactcheck{} only.  Classical transcendence inputs are labeled \knownstatus{}.}
\end{statusbox}

\subsection{Why this is a substantive narrowing rather than another reformulation}

The fixed-cross-ratio identity changes the mask problem at a practical level.  Previously a
tied Newton face was an open-ended local cancellation object whose coefficients could in
principle depend on a long moving valuation profile.  Here the moving profile is quotiented
out: along the primitive edge the two near-units contribute only a fixed rational $C_\ell$
(or its reciprocal).  Therefore the entire first local cancellation layer is governed by the
single integer polynomial $P(C_\ell^{\pm1})$.

There are then only two possibilities.  If this rational value is nonzero, its $\ell$-adic
valuation is bounded by the ordinary coefficient height and the edge length.  If it vanishes,
the corresponding binomial is an exact algebraic relation at the chosen near-unit pair and
can be divided out.  The one-prime geometric mechanism has no third branch.  This is stronger
than observing that a face must be long: it gives a normal form after the long edge has been
paid for.

The obstruction reappears only when different primes demand different tangent binomials.
Those directions are almost parallel, with lengths of order $q_k$, but their determinant is
a fixed nonzero integer.  The geometry is a thin constant-covolume lattice, not a large-area
lattice.  Simultaneous reduction can therefore collapse exponent support to finitely many
residue classes, but the basis coefficients needed for that collapse are of order $q_k$.
That coefficient growth is exactly on the same scale as the rational heights one wants to
remove.  The report ends with a precise cross-prime target at this interface.

\section{Counterexample normalization and synchronized near-units}
\label{sec:setup}

\subsection{Elementary reductions}

Assume for the rest of the structural development that Conjecture~\ref{conj:AE} fails, and
let $\beta>0$ be a nonintegral solution.  Put
\begin{equation}
 M=2^\beta\in\N,\qquad A=3^\beta\in\N,
 \qquad \thetaq=\frac{\log3}{\log2}.
 \label{eq:MA}
\end{equation}
The letters $M,A$ are kept in the order ``base two output, base three output.''

\begin{lemma}[Rational exponents are integral]
\label{lem:rational-integral}
If $x\in\Q$ and $2^x\in\Z$, then $x\in\Z$.
\end{lemma}

\begin{proof}
Because $2^x$ is a positive integer, $2^x\ge1$ and therefore $x\ge0$.  Write
$x=a/b$ in lowest terms with $a\ge0$ and $b>0$.  If $2^{a/b}=m\in\Z$, then
$m^b=2^a$.  Unique factorization gives $m=2^c$ and $a=bc$, so $b\mid a$.
Coprimality forces $b=1$.
\end{proof}

Thus a counterexample is irrational.  In fact it is transcendental: if $\beta$ were algebraic
irrational, Gelfond--Schneider would make $2^\beta$ transcendental, contrary to
$M\in\N$ \cite{Baker1975,LangTranscendental}.  This classical fact is background only; the
new transport arguments below use merely irrationality.

\begin{lemma}[Multiplicative independence]
\label{lem:MA-independent}
The integers $M$ and $A$ are multiplicatively independent.
\end{lemma}

\begin{proof}
If $M^rA^s=1$ for $r,s\in\Z$, then
$2^{r\beta}3^{s\beta}=1$, hence $2^r3^s=1$ after raising to the positive power $1/\beta$.
Unique factorization gives $r=s=0$.
\end{proof}

\subsection{Convergents and the rational near-unit orbit}

Let $p_k/q_k$ be the ordinary simple continued-fraction convergents of $\thetaq$, with
$q_k>0$, and put
\begin{equation}
 \eps_k=p_{k+1}q_k-p_kq_{k+1}\in\{+1,-1\}.
 \label{eq:epsilon}
\end{equation}
The standard determinant and approximation identities give
\begin{equation}
 |q_k\thetaq-p_k|<\frac1{q_{k+1}},
 \qquad
 \det\begin{pmatrix}p_k&p_{k+1}\\q_k&q_{k+1}\end{pmatrix}
 =p_kq_{k+1}-p_{k+1}q_k=-\eps_k.
 \label{eq:cf-basic}
\end{equation}

Define
\begin{equation}
 z_k=\frac{A^{q_k}}{M^{p_k}}.
 \label{eq:zk}
\end{equation}
Using \eqref{eq:MA},
\begin{equation}
 z_k
 =\exp\!\bigl(\beta(q_k\log3-p_k\log2)\bigr)
 =2^{\beta(q_k\thetaq-p_k)},
 \label{eq:zk-near}
\end{equation}
so $z_k\in\Q_{>0}$ and $z_k\to1$.  The small real parameter is
\begin{equation}
 \eta_k=q_k\thetaq-p_k,
 \qquad |\eta_k|<q_{k+1}^{-1}.
 \label{eq:eta}
\end{equation}

The key point is that the rational height of $z_k$ is linear in $q_k$, while
$|z_k-1|\asymp_\beta|\eta_k|$ is only polynomially small.  This familiar mismatch is why
ordinary rational approximation does not close the conjecture.  The goal here is different:
we determine the exact geometry of the primewise denominator depths of consecutive $z_k$.

\subsection{Local valuation vectors}

For every prime $\ell\mid MA$, write
\begin{equation}
 m_\ell=v_\ell(M),\qquad n_\ell=v_\ell(A),
 \qquad
 u_{\ell,k}=v_\ell(z_k)=q_kn_\ell-p_km_\ell.
 \label{eq:local-depth}
\end{equation}
The fixed vector $w_\ell=(m_\ell,n_\ell)\in\Nzero^2\setminus\{(0,0)\}$ records how the
prime $\ell$ occurs in the two outputs.  The moving pair
$(u_{\ell,k},u_{\ell,k+1})$ records its signed depth in two adjacent near-units.

The relation
\begin{equation}
 \sum_{\ell\mid MA}(n_\ell-\thetaq m_\ell)\log\ell
 =\log A-\thetaq\log M=0
 \label{eq:weighted-balance}
\end{equation}
will split the support into stable numerator and denominator primes.

\section{Unimodular transport of every local depth profile}
\label{sec:unimodular}

\subsection{The transport matrix}

\begin{theorem}[Exact unimodular transport]
\label{thm:unimodular-transport}
For every support prime $\ell$ and every $k$,
\begin{equation}
 \bigl(u_{\ell,k},u_{\ell,k+1}\bigr)
 =\bigl(m_\ell,n_\ell\bigr)B_k,
 \qquad
 B_k=
 \begin{pmatrix}
 -p_k&-p_{k+1}\\
 q_k&q_{k+1}
 \end{pmatrix}.
 \label{eq:transport-matrix}
\end{equation}
Moreover $\det B_k=\eps_k$, so $B_k\in\GL_2(\Z)$.
\end{theorem}

\begin{proof}
Multiplication gives
\[
 (m_\ell,n_\ell)B_k
 =(-p_km_\ell+q_kn_\ell,
   -p_{k+1}m_\ell+q_{k+1}n_\ell),
\]
which is the claimed depth pair.  The determinant is
\[
 (-p_k)q_{k+1}-(-p_{k+1})q_k
 =p_{k+1}q_k-p_kq_{k+1}=\eps_k.
\]
Since $\eps_k=\pm1$, the matrix is unimodular.
\end{proof}

\begin{statusbox}{\provedhere.  The theorem is elementary, but its consequences below do
not appear in this exact form in the supplied report.  In particular, the gcd is preserved
\emph{exactly}, not merely bounded by a fixed support exponent.}
\end{statusbox}

\subsection{Exact gcd and Smith invariance}

\begin{corollary}[Consecutive local gcd is constant]
\label{cor:gcd-invariance}
For every $\ell\mid MA$ and every $k$,
\begin{equation}
 \gcd\bigl(u_{\ell,k},u_{\ell,k+1}\bigr)
 =\gcd(m_\ell,n_\ell)=:h_\ell.
 \label{eq:gcd-invariance}
\end{equation}
The equality is understood with nonnegative gcds and is unchanged if both local depths are
replaced by their negatives.
\end{corollary}

\begin{proof}
A unimodular integral change of coordinates preserves the ideal generated by the two
coordinates.  Explicitly, each $u_{\ell,j}$ is an integral combination of $m_\ell,n_\ell$
by Theorem~\ref{thm:unimodular-transport}, while $(m_\ell,n_\ell)$ is an integral
combination of $(u_{\ell,k},u_{\ell,k+1})$ using $B_k^{-1}\in M_2(\Z)$.
The two generated ideals in $\Z$ are therefore equal.
\end{proof}

The same observation works simultaneously for every prime.  Let $S=\supp(MA)$ and form the
$|S|\times2$ valuation matrix
\begin{equation}
 V=\bigl((m_\ell,n_\ell)\bigr)_{\ell\in S},
 \qquad
 U_k=VB_k=\bigl((u_{\ell,k},u_{\ell,k+1})\bigr)_{\ell\in S}.
 \label{eq:VU}
\end{equation}

\begin{corollary}[Full determinantal-divisor invariance]
\label{cor:smith-invariance}
The matrices $V$ and $U_k$ have the same Smith normal form on the right.  In particular:
\begin{enumerate}
\item the gcd of all entries is independent of $k$;
\item the gcd of all $2\times2$ minors is independent of $k$;
\item the rank over every field, including every $\mathbb F_p$, is independent of $k$.
\end{enumerate}
\end{corollary}

\begin{proof}
Right multiplication by $B_k\in\GL_2(\Z)$ is an allowed unimodular column operation.  Smith
normal form and all determinantal ideals are invariant under such operations.
\end{proof}

This is a useful audit principle.  Passing from the output valuation vectors to adjacent
near-unit depths creates large coordinates but no new lattice index, no new local rank, and
no new Smith divisor.  Any method that uses only these invariants after the change of basis
has merely transported existing information.

\subsection{Cross-prime determinants}

\begin{theorem}[Cross-prime determinant is constant]
\label{thm:cross-determinant}
For any two support primes $\ell,r$,
\begin{equation}
 u_{\ell,k}u_{r,k+1}-u_{\ell,k+1}u_{r,k}
 =\eps_k\bigl(m_\ell n_r-n_\ell m_r\bigr).
 \label{eq:cross-determinant}
\end{equation}
\end{theorem}

\begin{proof}
Take determinants of the two-row submatrix of $U_k=VB_k$.  The determinant is the
corresponding minor of $V$ multiplied by $\det B_k=\eps_k$.
\end{proof}

Thus two local depth directions are proportional at one adjacent pair if and only if the
fixed valuation vectors $(m_\ell,n_\ell)$ and $(m_r,n_r)$ are proportional.  No accidental
rank drop can appear at a later convergent.

\subsection{The elementary transport identities recovered}

Expanding Theorem~\ref{thm:unimodular-transport} against the inverse continued-fraction
matrix gives
\begin{align}
 q_{k+1}u_{\ell,k}-q_ku_{\ell,k+1}&=\eps_km_\ell,
 \label{eq:q-transport}\\
 p_{k+1}u_{\ell,k}-p_ku_{\ell,k+1}&=\eps_kn_\ell.
 \label{eq:p-transport}
\end{align}
These identities are present in the source corpus in a transport formulation.  The new point
is to retain the full unimodular matrix, which immediately yields gcd, all-minor, and normal-
form consequences rather than using the identities one at a time.

\section{Stable signs and primitive local tangent vectors}
\label{sec:stable}

\subsection{The stable sign partition}

Define
\begin{equation}
 \delta_\ell=n_\ell-\thetaq m_\ell.
 \label{eq:delta}
\end{equation}
Because $\thetaq$ is irrational, $\delta_\ell\ne0$ for every nonzero integer vector
$(m_\ell,n_\ell)$.  Also
\begin{equation}
 u_{\ell,k}
 =q_k\delta_\ell+m_\ell\eta_k.
 \label{eq:u-asymptotic}
\end{equation}
Hence the sign of $u_{\ell,k}$ is eventually the sign of $\delta_\ell$.

Let
\begin{equation}
 \mathcal P_+=\{\ell:\delta_\ell>0\},
 \qquad
 \mathcal P_-=\{\ell:\delta_\ell<0\}.
 \label{eq:Ppm}
\end{equation}
For large $k$, primes in $\mathcal P_+$ occur in the reduced numerator of $z_k$, while primes
in $\mathcal P_-$ occur in its reduced denominator.

\begin{lemma}[Both stable sign classes occur]
\label{lem:both-signs}
The sets $\mathcal P_+$ and $\mathcal P_-$ are both nonempty.
\end{lemma}

\begin{proof}
Equation~\eqref{eq:weighted-balance} is a positive-logarithm weighted sum of the nonzero
numbers $\delta_\ell$ and equals zero.  They cannot all have the same sign.
\end{proof}

\subsection{Denominator depths}

Fix $\ell\in\mathcal P_-$ and put
\begin{equation}
 \gamma_\ell=\thetaq m_\ell-n_\ell=-\delta_\ell>0.
 \label{eq:gamma}
\end{equation}
For all sufficiently large $k$, define the positive denominator depth
\begin{equation}
 s_{\ell,k}=-u_{\ell,k}=p_km_\ell-q_kn_\ell.
 \label{eq:s-depth}
\end{equation}
Then
\begin{equation}
 s_{\ell,k}=q_k\gamma_\ell-m_\ell\eta_k,
 \qquad
 \left|s_{\ell,k}-q_k\gamma_\ell\right|<\frac{m_\ell}{q_{k+1}}.
 \label{eq:s-asymptotic}
\end{equation}

By Corollary~\ref{cor:gcd-invariance},
\begin{equation}
 \gcd(s_{\ell,k},s_{\ell,k+1})=h_\ell,
 \qquad h_\ell=\gcd(m_\ell,n_\ell).
 \label{eq:s-gcd}
\end{equation}
Define the primitive tangent coordinates
\begin{equation}
 a_{\ell,k}=\frac{s_{\ell,k+1}}{h_\ell},
 \qquad
 b_{\ell,k}=\frac{s_{\ell,k}}{h_\ell}.
 \label{eq:ab}
\end{equation}
They are positive coprime integers for all sufficiently large $k$.

\begin{proposition}[Linear tangent length]
\label{prop:tangent-length}
For every $\ell\in\mathcal P_-$,
\begin{align}
 a_{\ell,k}
 &=\frac{\gamma_\ell}{h_\ell}q_{k+1}
   +O_\ell(q_{k+2}^{-1}),
 \label{eq:a-asymp}\\
 b_{\ell,k}
 &=\frac{\gamma_\ell}{h_\ell}q_k
   +O_\ell(q_{k+1}^{-1}).
 \label{eq:b-asymp}
\end{align}
In particular, for all sufficiently large $k$,
\begin{equation}
 a_{\ell,k}+b_{\ell,k}
 \ge \frac{\gamma_\ell}{2h_\ell}(q_k+q_{k+1}).
 \label{eq:length-lower}
\end{equation}
\end{proposition}

\begin{proof}
Divide \eqref{eq:s-asymptotic} and its $k+1$ version by $h_\ell$.  The displayed lower
bound follows because the error is $o(1)$ while $q_k+q_{k+1}\to\infty$.
\end{proof}

The primitive local tangent is therefore not a bounded combinatorial direction.  It becomes
long at exactly the continued-fraction scale.

\section{Fixed rational cross-ratios along the moving tangents}
\label{sec:crossratio}

\subsection{The cross-ratio theorem}

Define the positive rational number
\begin{equation}
 C_\ell=\frac{A^{m_\ell/h_\ell}}{M^{n_\ell/h_\ell}}.
 \label{eq:Cell}
\end{equation}
This number depends only on the output pair and the primitive valuation direction at $\ell$;
it does not depend on $k$.

\begin{theorem}[Fixed tangent cross-ratio]
\label{thm:fixed-crossratio}
For every $\ell\in\mathcal P_-$ and every sufficiently large $k$,
\begin{equation}
 \boxed{
 \frac{z_k^{a_{\ell,k}}}{z_{k+1}^{b_{\ell,k}}}
 =C_\ell^{\eps_k}.}
 \label{eq:fixed-crossratio}
\end{equation}
Moreover
\begin{equation}
 v_\ell(C_\ell)=0,
 \qquad C_\ell>1.
 \label{eq:C-properties}
\end{equation}
\end{theorem}

\begin{proof}
Using $z_j=A^{q_j}/M^{p_j}$ and the definitions of $a,b$,
\[
 \frac{z_k^a}{z_{k+1}^b}
 =A^{(q_ks_{\ell,k+1}-q_{k+1}s_{\ell,k})/h_\ell}
  M^{(p_{k+1}s_{\ell,k}-p_ks_{\ell,k+1})/h_\ell}.
\]
The sign-reversed forms of \eqref{eq:q-transport}--\eqref{eq:p-transport} give
\[
 q_ks_{\ell,k+1}-q_{k+1}s_{\ell,k}=\eps_km_\ell,
 \qquad
 p_{k+1}s_{\ell,k}-p_ks_{\ell,k+1}=-\eps_kn_\ell.
\]
Substitution proves \eqref{eq:fixed-crossratio}.

For the local valuation,
\[
 v_\ell(C_\ell)
 =\frac{m_\ell n_\ell-n_\ell m_\ell}{h_\ell}=0.
\]
Finally,
\[
 \log C_\ell
 =\frac{m_\ell\log A-n_\ell\log M}{h_\ell}
 =\frac{\beta\log2}{h_\ell}(\thetaq m_\ell-n_\ell)
 =\frac{\beta\log2}{h_\ell}\gamma_\ell>0.
\]
Thus $C_\ell>1$.
\end{proof}

\begin{statusbox}{\provedhere.  This is the central new identity.  It turns the entire
moving $\ell$-adic tangent edge into evaluation at one fixed rational number.}
\end{statusbox}

\subsection{The numerator-prime analogue}

For $\ell\in\mathcal P_+$, the same theorem holds after setting
$s_{\ell,k}=u_{\ell,k}>0$ and reversing the ratio if desired.  Equivalently, define
$\gamma_\ell=|\thetaq m_\ell-n_\ell|$ and orient the primitive tangent so that both
coordinates are positive.  The denominator-prime orientation is the one relevant to reduced
denominator compression, so it will remain the default.

\subsection{The logarithmic meaning of the constants}

Equation~\eqref{eq:C-properties} also gives the exact logarithm
\begin{equation}
 \log C_\ell
 =\frac{\beta\log2}{h_\ell}\gamma_\ell.
 \label{eq:logC}
\end{equation}
Thus local cross-ratios are not new independent periods.  They are rational numbers whose
logarithms are rational-linear transforms of $\log M,\log A$, and their ratios are
fractional-linear transforms of $\thetaq$.

For two denominator primes $\ell,r$ define
\begin{equation}
 \rho_{\ell r}
 =\frac{\log C_r}{\log C_\ell}
 =\frac{h_\ell(\thetaq m_r-n_r)}
        {h_r(\thetaq m_\ell-n_\ell)}.
 \label{eq:rho}
\end{equation}
If $m_\ell n_r-n_\ell m_r\ne0$, this is a nonconstant $\PGL_2(\Q)$ transform of $\thetaq$.
Therefore any irrationality measure for $\rho_{\ell r}$ transports back to one for
$\thetaq$, and conversely.  Merely replacing $\log2/\log3$ by the ratio of two local
cross-ratio logarithms does not change the transcendence weight of the problem.

\subsection{Projectively transported convergents}

Put
\begin{equation}
 x_k=\frac{s_{r,k}}{h_r},
 \qquad y_k=\frac{s_{\ell,k}}{h_\ell}.
 \label{eq:xy-depth}
\end{equation}
Then
\begin{equation}
 \frac{x_k}{y_k}
 =\frac{h_\ell(m_rp_k-n_rq_k)}
        {h_r(m_\ell p_k-n_\ell q_k)}
 \longrightarrow \rho_{\ell r}.
 \label{eq:mobius-convergents}
\end{equation}
The approximation error has the exact form
\begin{equation}
 x_k-\rho_{\ell r}y_k
 =\frac{(m_\ell n_r-n_\ell m_r)(p_k-q_k\thetaq)}
        {h_r(\thetaq m_\ell-n_\ell)}.
 \label{eq:transported-error}
\end{equation}
Furthermore, Theorem~\ref{thm:cross-determinant} gives
\begin{equation}
 y_kx_{k+1}-y_{k+1}x_k
 =\eps_k\frac{m_\ell n_r-n_\ell m_r}{h_\ell h_r}.
 \label{eq:transported-det}
\end{equation}
Thus the local depth fractions are generalized convergents with a fixed cross-determinant.
When the normalized determinant has absolute value one, this is exactly a unimodular
continued-fraction transport.  The phenomenon is structurally useful, but it supplies no
new irrationality beyond the original convergents of $\thetaq$.

\section{Newton-face geometry and the linear-width threshold}
\label{sec:newton}

\subsection{Pole weights}

Let
\begin{equation}
 F(U,V)=\sum_{(i,j)\in\Omega}c_{ij}U^iV^j
 \in\Z[U,V],
 \qquad \Omega\subset\Nzero^2\text{ finite}.
 \label{eq:F}
\end{equation}
At a stable denominator prime $\ell$, the geometric pole weight of the monomial $U^iV^j$ is
\begin{equation}
 W_{\ell,k}(i,j)=i s_{\ell,k}+j s_{\ell,k+1}.
 \label{eq:pole-weight}
\end{equation}
Ignoring coefficient valuations for the moment, the exposed $\ell$-adic Newton face is the
set of exponents maximizing $W_{\ell,k}$ on $\Omega$.  A tie is possible only along the
kernel of this linear functional.

\begin{theorem}[Primitive tangent direction]
\label{thm:primitive-tangent}
For two exponent vectors $(i,j),(i',j')\in\Z^2$,
\begin{equation}
 W_{\ell,k}(i,j)=W_{\ell,k}(i',j')
 \label{eq:equal-weight}
\end{equation}
if and only if
\begin{equation}
 (i-i',j-j')=t\,(a_{\ell,k},-b_{\ell,k})
 \quad\text{for some }t\in\Z.
 \label{eq:tangent-direction}
\end{equation}
The vector $(a_{\ell,k},-b_{\ell,k})$ is primitive.
\end{theorem}

\begin{proof}
Equality of weights is
\[
 (i-i')s_{\ell,k}+(j-j')s_{\ell,k+1}=0.
\]
Divide by $h_\ell$ and use
$\gcd(s_{\ell,k}/h_\ell,s_{\ell,k+1}/h_\ell)=1$.  The complete integer solution set is
\[
 (i-i',j-j')
 =t\left(\frac{s_{\ell,k+1}}{h_\ell},
          -\frac{s_{\ell,k}}{h_\ell}\right)
 =t(a_{\ell,k},-b_{\ell,k}).
\]
Coprimality makes this direction primitive.
\end{proof}

\subsection{A sharp face-width obstruction}

\begin{corollary}[Every nondegenerate tied face is linearly wide]
\label{cor:face-width}
Suppose $F$ has bidegree at most $(d,e)$ and its $\ell$-adic exposed face contains two
distinct monomials.  Then
\begin{equation}
 d\ge a_{\ell,k},
 \qquad e\ge b_{\ell,k},
 \qquad d+e\ge a_{\ell,k}+b_{\ell,k}.
 \label{eq:face-width}
\end{equation}
Consequently, for all sufficiently large $k$,
\begin{equation}
 d+e\ge
 \frac{\gamma_\ell}{2h_\ell}(q_k+q_{k+1}).
 \label{eq:linear-degree-tax}
\end{equation}
\end{corollary}

\begin{proof}
By Theorem~\ref{thm:primitive-tangent}, the difference between two distinct face exponents
is a nonzero multiple of $(a,-b)$.  Its horizontal and vertical absolute displacements are
at least $a$ and $b$.  Both exponents lie in $[0,d]\times[0,e]$, giving the first two
inequalities.  Apply Proposition~\ref{prop:tangent-length} for the last one.
\end{proof}

\begin{corollary}[Sublinear-degree masks have no shared structural face]
\label{cor:sublinear-no-face}
Let $F_k\in\Z[U,V]$ have bidegrees $(d_k,e_k)$ with
\begin{equation}
 d_k+e_k=o(q_k).
 \label{eq:sublinear-degree}
\end{equation}
Then, for every stable denominator prime $\ell$, the $\ell$-adic exposed geometric face of
$F_k$ is a single monomial for all sufficiently large $k$.
\end{corollary}

\begin{proof}
There are finitely many support primes.  Apply Corollary~\ref{cor:face-width} to each one.
\end{proof}

This closes the entire sublinear-degree branch of the shared-face target.  It also shows that
``fixed degree with increasingly deep cancellation'' is impossible at the geometric-face
level.  Any fixed-degree construction can gain local valuation only from divisibility of its
coefficients or from cancellation after coefficient valuations have altered the tropical
weights.

\subsection{Several prime slopes on one Newton boundary}

The degree tax adds across distinct slopes.  Let $\mathcal N(F)=\conv(\Omega)$ be the Newton
polygon.  For positive normal vectors $(s_{\ell,k},s_{\ell,k+1})$, all exposed faces occur on
the monotone southeast boundary: moving along that boundary, the first coordinate increases
and the second decreases.

\begin{theorem}[Multi-prime perimeter tax]
\label{thm:perimeter-tax}
Let $\ell_1,\ldots,\ell_R\in\mathcal P_-$ represent distinct valuation slopes, meaning that
the primitive vectors
$(m_{\ell_i}/h_{\ell_i},n_{\ell_i}/h_{\ell_i})$ are pairwise distinct.  Suppose the Newton
polygon of $F$ has a nondegenerate exposed face for each of these primes.  If the face for
$\ell_i$ has lattice length $L_i\ge1$ in its primitive tangent direction, then
\begin{align}
 d&\ge\sum_{i=1}^R L_i a_{\ell_i,k},
 \label{eq:sum-a}\\
 e&\ge\sum_{i=1}^R L_i b_{\ell_i,k},
 \label{eq:sum-b}\\
 d+e&\ge\sum_{i=1}^R L_i(a_{\ell_i,k}+b_{\ell_i,k}).
 \label{eq:sum-perimeter}
\end{align}
In particular, each additional distinct tied prime slope consumes another linear-width edge.
\end{theorem}

\begin{proof}
Distinct primitive valuation slopes give distinct normal slopes and hence distinct exposed
edges.  Orient the relevant Newton boundary from northwest to southeast.  The edge for
$\ell_i$ has displacement
$L_i(a_{\ell_i,k},-b_{\ell_i,k})$.  Horizontal displacements of all such edges are
nonnegative and sum to at most the total horizontal width $d$; the absolute vertical
displacements sum to at most the total vertical width $e$.  This proves
\eqref{eq:sum-a}--\eqref{eq:sum-b} and their sum.
\end{proof}

The theorem is sharp at the level of polygon geometry: two long primitive vectors can have
constant determinant and can bound a very thin lattice triangle.  Thus perimeter, not area,
is the correct first invariant.  The constant-area phenomenon becomes central in
Section~\ref{sec:crossprime}.

\section{Edge-polynomial reduction at one structural prime}
\label{sec:edgepoly}

\subsection{Evaluation of a tied edge}

Fix $\ell\in\mathcal P_-$ and $k$ sufficiently large.  Abbreviate
\begin{equation}
 a=a_{\ell,k},\qquad b=b_{\ell,k},
 \qquad C=C_\ell^{\eps_k}.
 \label{eq:abc-short}
\end{equation}
Let a tied edge have exponent points
\begin{equation}
 (i_t,j_t)=(i_0,j_0)+t(a,-b),
 \qquad 0\le t\le L,
 \label{eq:edge-points}
\end{equation}
where some coefficients may vanish.  Its Laurent edge polynomial is
\begin{equation}
 E(U,V)=\sum_{t=0}^L c_tU^{i_t}V^{j_t}.
 \label{eq:edge-E}
\end{equation}

\begin{theorem}[Exact edge-polynomial evaluation]
\label{thm:edge-evaluation}
With
\begin{equation}
 P(T)=\sum_{t=0}^L c_tT^t\in\Z[T],
 \label{eq:edge-P}
\end{equation}
one has
\begin{equation}
 E(z_k,z_{k+1})
 =z_k^{i_0}z_{k+1}^{j_0}P(C).
 \label{eq:edge-evaluation}
\end{equation}
\end{theorem}

\begin{proof}
Factor out the first monomial and use Theorem~\ref{thm:fixed-crossratio}:
\[
 U^{i_t}V^{j_t}\big|_{(z_k,z_{k+1})}
 =z_k^{i_0}z_{k+1}^{j_0}
   \left(\frac{z_k^a}{z_{k+1}^b}\right)^t
 =z_k^{i_0}z_{k+1}^{j_0}C^t.
\]
Summing gives the claim.
\end{proof}

This is the decisive simplification.  The moving edge length and moving exponents disappear
from the local initial form; only the fixed rational $C_\ell^{\pm1}$ remains.

\subsection{A rational valuation bound}

For a nonzero rational $R=N/D$ in lowest terms, put
\begin{equation}
 H(R)=\max(|N|,|D|).
 \label{eq:mult-height}
\end{equation}

\begin{lemma}[Valuation of a polynomial at an $\ell$-unit rational]
\label{lem:rational-eval-bound}
Let $R=N/D\in\Q^\times$ be in lowest terms and suppose $\ell\nmid ND$.  If
$P(T)=\sum_{t=0}^Lc_tT^t\in\Z[T]$ and $P(R)\ne0$, then
\begin{equation}
 v_\ell(P(R))\log\ell
 \le \log\lVert P\rVert_1+L\log H(R),
 \qquad \lVert P\rVert_1=\sum_{t=0}^L|c_t|.
 \label{eq:rational-eval-bound}
\end{equation}
\end{lemma}

\begin{proof}
Write
\[
 P(N/D)=D^{-L}J,
 \qquad
 J=\sum_{t=0}^Lc_tN^tD^{L-t}\in\Z\setminus\{0\}.
\]
Because $\ell\nmid D$, $v_\ell(P(R))=v_\ell(J)$.  Also
\[
 \ell^{v_\ell(J)}\le |J|
 \le\lVert P\rVert_1H(R)^L.
\]
Take logarithms.
\end{proof}

Since $v_\ell(C_\ell)=0$, this applies to $C=C_\ell^{\eps_k}$ with the same height
$H(C_\ell)$.  Thus a nonzero edge initial form cannot acquire arbitrarily many local layers
for free.

\begin{corollary}[One-prime edge dichotomy]
\label{cor:edge-dichotomy}
For a tied $\ell$-adic edge of lattice length $L$ exactly one of the following holds.
\begin{enumerate}
\item The edge value is nonzero, and its extra $\ell$-adic valuation beyond the common
monomial factor satisfies
\begin{equation}
 v_\ell(P(C_\ell^{\eps_k}))\log\ell
 \le \log\lVert P\rVert_1+L\log H(C_\ell).
 \label{eq:edge-extra-bound}
\end{equation}
\item The edge value vanishes exactly, and the primitive rational-root factor associated to
$C_\ell^{\eps_k}$ divides $P$ in $\Z[T]$.
\end{enumerate}
\end{corollary}

\begin{proof}
The nonzero branch is Lemma~\ref{lem:rational-eval-bound}.  In the zero branch write
$C_\ell^{\eps_k}=N/D$ in lowest terms.  Gauss's lemma and the rational-root theorem imply
$DT-N\mid P(T)$ in $\Z[T]$.
\end{proof}

\subsection{The edge-length gain is subcritical}

A face contained in a $(d,e)$ box has
\begin{equation}
 L\le\min\left(\left\lfloor\frac d a\right\rfloor,
               \left\lfloor\frac e b\right\rfloor\right)
 \le\frac{d+e}{a+b}.
 \label{eq:edge-length-bound}
\end{equation}
Using Proposition~\ref{prop:tangent-length},
\begin{equation}
 L=O_\ell\!\left(\frac{d+e}{q_k+q_{k+1}}\right).
 \label{eq:L-asymp}
\end{equation}
Therefore the nonzero edge branch contributes at most
\begin{equation}
 \log\lVert P\rVert_1
 +O_\ell\!\left(\frac{d+e}{q_k}\right)
 \label{eq:edge-small-gain}
\end{equation}
in logarithmic $\ell$-adic gain.  The pole depth one seeks to cancel is typically of order
$q_k(d+e)$.  Unless the coefficients themselves carry critical height, nonzero one-prime
edge evaluation is negligible by two powers of the convergent scale.

This estimate is intentionally local.  It does not forbid a coefficient system engineered to
carry large $\ell$-adic divisibility, nor does it control simultaneous cancellation at other
primes.  It says that the geometry of a single tied face, by itself, contributes no hidden
large factor.

\section{The tangent binomial and a canonical one-prime normal form}
\label{sec:normalform}

\subsection{The exact tangent binomial}

Again write $C_\ell^{\eps_k}=N/D$ in lowest positive terms.  The fixed-cross-ratio theorem is
equivalent to
\begin{equation}
 D z_k^a=N z_{k+1}^b.
 \label{eq:tangent-relation}
\end{equation}
Define the primitive tangent binomial
\begin{equation}
 B_{\ell,k}(U,V)=DU^a-NV^b\in\Z[U,V].
 \label{eq:tangent-binomial}
\end{equation}
It vanishes exactly at $(z_k,z_{k+1})$.  Because $\gcd(N,D)=1$ and $\gcd(a,b)=1$, it is a
primitive irreducible binomial over $\Q$.

The zero branch of Corollary~\ref{cor:edge-dichotomy} is precisely divisibility of the edge
initial form by this binomial, after factoring its endpoint monomial.  The next theorem shows
that the binomial can be removed globally, not merely on one face.

\subsection{Integral normal-form theorem}

\begin{theorem}[Canonical one-prime tangent normal form]
\label{thm:one-prime-normal-form}
Let
\[
 F(U,V)=\sum_{0\le i\le d,\ 0\le j\le e}c_{ij}U^iV^j\in\Z[U,V],
\]
and let $a,b,N,D$ be as above.  Put $T=\lfloor d/a\rfloor$.  For each $i$, write
$i=t(i)a+r(i)$ with $0\le r(i)<a$, and define
\begin{equation}
 G(U,V)
 =\sum_{i,j}c_{ij}N^{t(i)}D^{T-t(i)}
   U^{r(i)}V^{j+b t(i)},
 \label{eq:normal-G-raw}
\end{equation}
combining like monomials.  Then:
\begin{enumerate}
\item $G\in\Z[U,V]$ and
\begin{equation}
 G(z_k,z_{k+1})=D^T F(z_k,z_{k+1});
 \label{eq:normal-eval}
\end{equation}
\item $\deg_U G<a$ and $\deg_VG\le e+bT$;
\item
\begin{equation}
 \lVert G\rVert_1\le H(C_\ell)^T\lVert F\rVert_1;
 \label{eq:normal-height}
\end{equation}
\item no two distinct monomials of $G$ have the same geometric $\ell$-adic pole weight;
\item $G=0$ if and only if $B_{\ell,k}$ divides $F$ in $\Q[U,V]$; because
$B_{\ell,k}$ is primitive, this is equivalent to divisibility in $\Z[U,V]$ up to the
content of $F$;
\item since $\ell\nmid D$,
\begin{equation}
 v_\ell(G(z_k,z_{k+1}))=v_\ell(F(z_k,z_{k+1})).
 \label{eq:normal-v}
\end{equation}
\end{enumerate}
\end{theorem}

\begin{proof}
At the evaluation point, relation \eqref{eq:tangent-relation} gives
\[
 z_k^{i}z_{k+1}^{j}
 =z_k^{r+ta}z_{k+1}^j
 =\left(\frac ND\right)^t z_k^r z_{k+1}^{j+bt}.
\]
Multiplying by $D^T$ and summing proves \eqref{eq:normal-eval} and the integrality of
\eqref{eq:normal-G-raw}.  The degree bounds are immediate.  Every coefficient multiplier
$N^tD^{T-t}$ has absolute value at most $H(C_\ell)^T$, proving
\eqref{eq:normal-height} after the triangle inequality.

For the pole-weight claim, suppose two exponents $(r,j)$ and $(r',j')$ of $G$ have equal
weight.  Theorem~\ref{thm:primitive-tangent} gives
$(r-r',j-j')=m(a,-b)$.  But $0\le r,r'<a$, so $|r-r'|<a$ and therefore $m=0$.  Hence the
exponents are equal.

Over $\Q[V]$, division by the polynomial
$U^a-(N/D)V^b$ has a unique remainder of $U$-degree $<a$.  Up to the nonzero scalar $D^T$,
$G$ is exactly that remainder.  Thus $G=0$ if and only if the binomial divides $F$ over
$\Q[U,V]$.  Gauss's lemma handles integral divisibility.  Finally
$v_\ell(D)=0$ by Theorem~\ref{thm:fixed-crossratio}, so \eqref{eq:normal-eval} preserves the
$\ell$-adic valuation.
\end{proof}

\begin{statusbox}{\provedhere.  This theorem is the main algebraic output: exact geometric
cancellation along one structural-prime tangent can always be divided out, with an explicit
and comparatively small coefficient-height tax.}
\end{statusbox}

\subsection{Asymptotic cost of the reduction}

By Proposition~\ref{prop:tangent-length}, $a\asymp_\ell q_{k+1}$.  Hence
\begin{equation}
 T=\left\lfloor\frac d a\right\rfloor
 =O_\ell\!\left(\frac d{q_{k+1}}\right),
 \label{eq:T-cost}
\end{equation}
and the logarithmic coefficient inflation is
\begin{equation}
 \log\lVert G\rVert_1-
 \log\lVert F\rVert_1
 \le O_\ell\!\left(\frac d{q_{k+1}}\right)\log H(C_\ell).
 \label{eq:normal-cost}
\end{equation}
This is tiny compared with a structural denominator scale of order $q_kd$, provided the
reduction is performed at only one prime.

The vertical degree increases by at most $bT$.  Since $b/a\asymp q_k/q_{k+1}<1$, this is of
order $d$ rather than $q_kd$.  The normal form therefore does not hide an exponential support
explosion.

\subsection{What the normal form proves and what it does not}

\begin{corollary}[One-prime geometric capture is quotient-trivial]
\label{cor:one-prime-trivial}
Fix a stable denominator prime $\ell$.  Modulo the exact relation
$B_{\ell,k}(z_k,z_{k+1})=0$, every polynomial mask has a representative with no tied
$\ell$-adic geometric pole face, with the same local value up to an $\ell$-adic unit and with
coefficient-height cost given by \eqref{eq:normal-cost}.
\end{corollary}

This does \emph{not} imply that the local valuation of the reduced representative is small.
The coefficients may themselves be highly divisible by $\ell$, and distinct geometric
weights may become tied after their coefficient valuations are added.  Moreover the units
occurring in different monomials can satisfy deep congruences.  Those are genuine arithmetic
phenomena.  What has been removed is the idea that a long shared Newton face supplies deep
cancellation merely because its monomials have equal pole order.

A future one-prime argument must therefore prove a coefficient statement.  For example, it
could show that high $\ell$-adic valuation of the normal-form value forces a comparable
$\ell$-adic content in all leading coefficients, allowing primitive descent.  No such theorem
is presently known.  The next section shows why performing normal-form reduction at several
primes is subtler: the relevant tangent bases are nearly parallel.

\section{Cross-prime transversality and the constant-index tangent lattice}
\label{sec:crossprime}

\subsection{Primitive tangent vectors}

For each stable denominator prime define
\begin{equation}
 t_{\ell,k}=(a_{\ell,k},-b_{\ell,k})\in\Z^2.
 \label{eq:tangent-vector}
\end{equation}
This is the primitive exponent direction on which the $\ell$-adic pole weight is constant.
It has Euclidean and $\ell^1$ length $\asymp_\ell q_k+q_{k+1}$.

Let $W_-$ be the matrix whose rows are the primitive valuation vectors
$(m_\ell/h_\ell,n_\ell/h_\ell)$ for $\ell\in\mathcal P_-$, and let $T_k$ be the matrix whose
corresponding rows are $t_{\ell,k}$.  Put
\[
 J=\begin{pmatrix}0&1\\-1&0\end{pmatrix}.
\]

\begin{theorem}[Full tangent-matrix equivalence]
\label{thm:tangent-matrix-equivalence}
For every sufficiently large $k$,
\begin{equation}
 T_k=W_-B_kJ,
 \label{eq:tangent-matrix-equivalence}
\end{equation}
where $B_k$ is the unimodular transport matrix of
Theorem~\ref{thm:unimodular-transport}.  Consequently $T_k$ and $W_-$ have the same Smith
normal form and the same determinantal divisors.
\end{theorem}

\begin{proof}
For one row, Theorem~\ref{thm:unimodular-transport} gives
\[
 \left(\frac{m_\ell}{h_\ell},\frac{n_\ell}{h_\ell}\right)B_k
 =\left(\frac{u_{\ell,k}}{h_\ell},\frac{u_{\ell,k+1}}{h_\ell}\right)
 =(-b_{\ell,k},-a_{\ell,k}).
\]
Right multiplication by $J$ sends $(-b,-a)$ to $(a,-b)$.  Both $B_k$ and $J$ are
unimodular, so the Smith and determinantal claims follow.
\end{proof}

This theorem says that the entire moving tangent configuration---not only each pairwise
determinant---is a modular image of the fixed primitive valuation configuration.  Long,
almost parallel tangent vectors can therefore create no new lattice invariant whatsoever.

\begin{theorem}[Constant cross-prime tangent determinant]
\label{thm:tangent-determinant}
For two stable denominator primes $\ell,r$,
\begin{equation}
 \det(t_{\ell,k},t_{r,k})
 =\eps_k\,
 \frac{m_\ell n_r-n_\ell m_r}{h_\ell h_r}.
 \label{eq:tangent-determinant}
\end{equation}
In particular, its absolute value is independent of $k$.
\end{theorem}

\begin{proof}
Using $t_{\ell,k}=(s_{\ell,k+1}/h_\ell,-s_{\ell,k}/h_\ell)$,
\begin{align*}
 \det(t_{\ell,k},t_{r,k})
 &=\frac{s_{\ell,k}s_{r,k+1}-s_{\ell,k+1}s_{r,k}}
        {h_\ell h_r}\\
 &=\frac{u_{\ell,k}u_{r,k+1}-u_{\ell,k+1}u_{r,k}}
        {h_\ell h_r}.
\end{align*}
Apply Theorem~\ref{thm:cross-determinant}.
\end{proof}

Define the normalized slope determinant
\begin{equation}
 \Delta_{\ell r}
 =\frac{m_\ell n_r-n_\ell m_r}{h_\ell h_r}\in\Z.
 \label{eq:Delta-lr}
\end{equation}
The integrality follows because $h_\ell$ and $h_r$ divide the corresponding rows.  The
primitive valuation directions are distinct exactly when $\Delta_{\ell r}\ne0$.

\begin{corollary}[Constant-index tangent lattice]
\label{cor:constant-index}
If $\Delta_{\ell r}\ne0$, the lattice
\begin{equation}
 \Lambda_{\ell r,k}
 =\Z t_{\ell,k}+\Z t_{r,k}\subset\Z^2
 \label{eq:tangent-lattice}
\end{equation}
has index
\begin{equation}
 [\Z^2:\Lambda_{\ell r,k}]=|\Delta_{\ell r}|,
 \label{eq:lattice-index}
\end{equation}
independent of $k$.
\end{corollary}

This is the central cross-prime geometry: two tangent generators become arbitrarily long and
almost parallel, yet they span a lattice of fixed covolume.  Their angle is of order
$q_k^{-2}$ when $q_{k+1}\asymp q_k$.

\subsection{Cross-ratio dependence and valuation-slope dependence coincide}

\begin{proposition}[Equivalent slope classifications]
\label{prop:slope-classification}
For two support primes $\ell,r$, the following are equivalent:
\begin{enumerate}
\item $(m_\ell,n_\ell)$ and $(m_r,n_r)$ are proportional over $\Q$;
\item their primitive vectors
$(m_\ell/h_\ell,n_\ell/h_\ell)$ and
$(m_r/h_r,n_r/h_r)$ are equal;
\item $\Delta_{\ell r}=0$;
\item $t_{\ell,k}$ and $t_{r,k}$ are proportional for one, hence every, sufficiently large
$k$;
\item $C_\ell$ and $C_r$ are multiplicatively dependent.
\end{enumerate}
\end{proposition}

\begin{proof}
The equivalence of the first four statements is immediate from primitive integer vectors and
Theorem~\ref{thm:tangent-determinant}.  For the last statement, write
\[
 C_\ell=A^{m_\ell/h_\ell}M^{-n_\ell/h_\ell},
 \qquad
 C_r=A^{m_r/h_r}M^{-n_r/h_r}.
\]
By Lemma~\ref{lem:MA-independent}, an integer relation between powers of $C_\ell,C_r$ is
exactly an integer relation between the two primitive exponent vectors.  Such a relation
exists if and only if they are proportional.
\end{proof}

Thus prime slopes, tangent directions, and cross-ratio multiplicative classes are the same
finite partition of the support.  Counting several primes in one class does not provide
several independent local equations.

\subsection{Exact short-vector identities}

Fix two nonproportional denominator primes and abbreviate
\begin{equation}
 t_\ell=(a_\ell,-b_\ell),\qquad
 t_r=(a_r,-b_r),\qquad
 \Delta=\det(t_\ell,t_r)=\eps_k\Delta_{\ell r}.
 \label{eq:short-abbrev}
\end{equation}
The two long vectors satisfy the exact Bezout identities
\begin{align}
 -b_r t_\ell+b_\ell t_r&=\Delta(1,0),
 \label{eq:short-e1}\\
 -a_r t_\ell+a_\ell t_r&=\Delta(0,1).
 \label{eq:short-e2}
\end{align}
These are direct coordinate calculations.  They show that a bounded vector in the
constant-index lattice requires coefficients of size $\asymp q_k$ in the long tangent basis.

Apply the corresponding monomial relations
\begin{equation}
 z_k^{a_\ell}z_{k+1}^{-b_\ell}=C_\ell^{\eps_k},
 \qquad
 z_k^{a_r}z_{k+1}^{-b_r}=C_r^{\eps_k}.
 \label{eq:two-monomial-relations}
\end{equation}

\begin{theorem}[Exact two-prime resultant identities]
\label{thm:resultant-identities}
With the signed determinant $\Delta$ above,
\begin{align}
 z_k^{\Delta}
 &=\left(C_\ell^{-b_r}C_r^{b_\ell}\right)^{\eps_k},
 \label{eq:resultant-zk}\\
 z_{k+1}^{\Delta}
 &=\left(C_\ell^{-a_r}C_r^{a_\ell}\right)^{\eps_k}.
 \label{eq:resultant-zkp1}
\end{align}
\end{theorem}

\begin{proof}
Raise the first relation in \eqref{eq:two-monomial-relations} to $-b_r$ and the second to
$b_\ell$, then multiply.  The exponent of $z_{k+1}$ cancels and the exponent of $z_k$ is
$-a_\ell b_r+a_rb_\ell=\Delta$.  This proves \eqref{eq:resultant-zk}.  The second identity
uses exponents $-a_r$ and $a_\ell$.
\end{proof}

These are binomial resultants in closed form.  They are also a conservation law: eliminating
one near-unit from two fixed-cross-ratio equations recovers a bounded power of the other
near-unit, but the eliminating exponents are themselves of size $q_k$.

\subsection{Critical conditioning}

Let
\begin{equation}
 T_k=
 \begin{pmatrix}
 a_\ell&a_r\\
 -b_\ell&-b_r
 \end{pmatrix}.
 \label{eq:tangent-matrix}
\end{equation}
Its entries are $\Theta(q_k+q_{k+1})$ while $|\det T_k|=|\Delta_{\ell r}|$ is fixed.
Consequently
\begin{equation}
 \lVert T_k^{-1}\rVert_\infty
 =\Theta_{\ell,r}\!\left(\frac{q_k+q_{k+1}}{|\Delta_{\ell r}|}\right).
 \label{eq:condition-number}
\end{equation}
This estimate follows directly from the adjugate formula; no numerical linear algebra is
involved.

The quotient of the Laurent exponent lattice by the two tangent relations has only
$|\Delta_{\ell r}|$ residue classes.  Algebraically, every Laurent monomial can be reduced
to one of finitely many residue monomials times a rational power product in $C_\ell,C_r$.
At first sight this looks much stronger than the one-prime normal form.  The inverse-matrix
estimate identifies the cost: reducing an exponent displacement of size $D$ can require
powers of $C_\ell,C_r$ of size $O(q_kD/|\Delta_{\ell r}|)$.

The short-vector identities show that this is not an artifact of a loose norm estimate.
Even the bounded displacement $\Delta(1,0)$ uses coefficients $-b_r,b_\ell$, each of order
$q_k$.  Raising \eqref{eq:short-e1} by $D$ gives a displacement of size $O(D)$ with tangent
coefficients of size $\Theta(q_kD)$.  Thus simultaneous normal-form reduction sits at the
same linear-height scale as the original rational near-units.

\begin{proposition}[Height conservation in the two-prime elimination]
\label{prop:height-conservation}
For the absolute logarithmic height $\height$ on $\Q^\times$,
\begin{align}
 \height\!\left(C_\ell^{-b_r}C_r^{b_\ell}\right)
 &=|\Delta|\,\height(z_k),
 \label{eq:height-zk}\\
 \height\!\left(C_\ell^{-a_r}C_r^{a_\ell}\right)
 &=|\Delta|\,\height(z_{k+1}).
 \label{eq:height-zkp1}
\end{align}
\end{proposition}

\begin{proof}
Theorem~\ref{thm:resultant-identities} identifies each rational number with
$z_j^{\pm\Delta}$.  The height on $\Q^\times$ satisfies
$\height(x^n)=|n|\height(x)$ and is unchanged by inversion.
\end{proof}

The proposition warns against estimating the right-hand sides by the triangle inequality in
the exponents of $C_\ell,C_r$: those large fixed-rational powers undergo exact cancellation
across places.  The reduced height is not quadratic in $q_k$; it returns to the critical
linear scale.  A proof must therefore exploit the \emph{distribution} of this cancellation,
not merely its total height.

\subsection{A thin unimodular cell}

When $|\Delta_{\ell r}|=1$, the tangent vectors form a basis of $\Z^2$.  The triangle with
vertices $0,t_{\ell,k},t_{r,k}$ has area $1/2$ even though two sides have length
$\asymp q_k$.  It is a unimodular triangle of arbitrarily bad aspect ratio.  This explains
why area-based Newton polygon estimates miss the true cost: the polygon can support two
prime-specific long edges with only constant area and three vertices.

The same phenomenon suggests a useful computational basis.  Rather than search arbitrary
large boxes, one should enumerate coefficient patterns on these thin tangent cells and their
small Minkowski sums.  Section~\ref{sec:program} formulates the corresponding exact integer
program.

\section{Three-prime syzygies and the rank-two ceiling}
\label{sec:syzygy}

\subsection{Primitive cross-ratio exponent vectors}

Define
\begin{equation}
 c_\ell=\left(\frac{m_\ell}{h_\ell},-\frac{n_\ell}{h_\ell}\right)\in\Z^2,
 \qquad C_\ell=A^{c_{\ell,1}}M^{c_{\ell,2}}.
 \label{eq:c-vector}
\end{equation}
The vectors $c_\ell$ live in the fixed rank-two exponent lattice of the multiplicatively
independent pair $(A,M)$.

For any two primes put
\begin{equation}
 [\ell,r]=\det(c_\ell,c_r).
 \label{eq:bracket}
\end{equation}
Up to a harmless sign, this is the normalized slope determinant
$\Delta_{\ell r}$.

\begin{theorem}[Pl\"ucker syzygy for three cross-ratios]
\label{thm:plucker}
For any three support primes $\ell,r,s$,
\begin{equation}
 [r,s]c_\ell+[s,\ell]c_r+[\ell,r]c_s=0.
 \label{eq:plucker-vector}
\end{equation}
Consequently
\begin{equation}
 C_\ell^{[r,s]}
 C_r^{[s,\ell]}
 C_s^{[\ell,r]}=1.
 \label{eq:plucker-C}
\end{equation}
\end{theorem}

\begin{proof}
Identity \eqref{eq:plucker-vector} is the standard alternating relation among three vectors
in a two-dimensional vector space; it follows by expanding both coordinates.  Apply the
homomorphism $(a,b)\mapsto A^aM^b$ to obtain \eqref{eq:plucker-C}.
\end{proof}

Thus adding more support primes does not produce an arbitrarily high-rank family of tangent
constants.  All $C_\ell$ lie in the rank-two group $\langle A,M\rangle$, and every triple has
an explicit determinantal multiplicative relation.

\subsection{Oriented-matroid interpretation}

The signs of the brackets $[\ell,r]$ encode the cyclic order of valuation slopes.  The same
sign pattern controls:
\begin{enumerate}
\item the order of stable prime normals on the Newton boundary;
\item the orientation of the long tangent edges;
\item which cross-ratio powers occur in the two-prime resultants;
\item the three-term multiplicative syzygies among the $C_\ell$.
\end{enumerate}
The local mask problem therefore has a rank-two oriented-matroid skeleton fixed entirely by
the output valuation vectors.  The convergent index changes the lengths of the tangent
vectors but not this combinatorial skeleton.

This observation limits what a ``many primes'' argument can gain from independence alone.
The number of primes may be large, but the tangent constants have rank two.  Any genuine
many-prime advantage must come from simultaneous congruence conditions, coefficient
primitivity, or place-by-place incompatibility, not from an increase in multiplicative rank.

\subsection{A finite presentation of the tangent quotient}

Choose two nonproportional prime classes $\ell,r$.  Every other $c_s$ is a rational-linear
combination of $c_\ell,c_r$, and the Pl\"ucker identity clears the denominators explicitly.
Consequently every tangent constant $C_s$ is a rational power product of $C_\ell,C_r$, with
exponents controlled by the fixed determinants $[\cdot,\cdot]$.  At the exponent-lattice
level, all tangent binomials are therefore generated after saturation by two of them.

The word ``after saturation'' matters.  Determinants $|\Delta_{\ell r}|>1$ create finite
index and root issues.  However all quantities here are already positive rationals, so the
required roots, when they occur, are visible in the valuation lattice.  A Smith-normal-form
implementation can track them exactly.  This makes the tangent quotient particularly
well-suited to formalization and exact computation.

\section{A synthetic exact calibration}
\label{sec:calibration}

This section checks the formulas on a small symbolic example.  It is not numerical evidence
for or against the conjecture; the valuation vectors are synthetic and no integers $M,A$
with these exact local data are asserted to form a counterexample.

Take two primitive denominator-slope vectors
\begin{equation}
 (m_\ell,n_\ell)=(2,3),
 \qquad
 (m_r,n_r)=(3,4).
 \label{eq:sample-vectors}
\end{equation}
Both have $n/m<\thetaq$.  Use the consecutive convergents
\begin{equation}
 \frac{p_k}{q_k}=\frac{485}{306},
 \qquad
 \frac{p_{k+1}}{q_{k+1}}=\frac{1054}{665},
 \qquad \eps_k=-1.
 \label{eq:sample-convergents}
\end{equation}
The depths and tangents are:

\begin{center}
\begin{tabular}{@{}cccccc@{}}
\toprule
vector $(m,n)$ & $h$ & $s_k=mp_k-nq_k$ & $s_{k+1}$ & $a=s_{k+1}/h$ & $b=s_k/h$\\
\midrule
$(2,3)$ & $1$ & $52$ & $113$ & $113$ & $52$\\
$(3,4)$ & $1$ & $231$ & $502$ & $502$ & $231$\\
$(5,7)$ & $1$ & $283$ & $615$ & $615$ & $283$\\
\bottomrule
\end{tabular}
\end{center}

Thus
\begin{equation}
 t_\ell=(113,-52),
 \qquad t_r=(502,-231),
 \qquad \det(t_\ell,t_r)=1.
 \label{eq:sample-tangents}
\end{equation}
Two vectors of lengths several hundred form a unimodular basis.

The fixed constants are
\begin{equation}
 C_\ell=\frac{A^2}{M^3},
 \qquad C_r=\frac{A^3}{M^4}.
 \label{eq:sample-C}
\end{equation}
The cross-ratio identities read
\begin{equation}
 \frac{z_k^{113}}{z_{k+1}^{52}}=C_\ell^{-1},
 \qquad
 \frac{z_k^{502}}{z_{k+1}^{231}}=C_r^{-1}.
 \label{eq:sample-cross}
\end{equation}
The exact short-vector identity
\begin{equation}
 -231(113,-52)+52(502,-231)=(1,0)
 \label{eq:sample-bezout}
\end{equation}
then gives
\begin{equation}
 z_k=C_\ell^{231}C_r^{-52}.
 \label{eq:sample-z}
\end{equation}
Indeed the right side simplifies to
\[
 \frac{A^{462}}{M^{693}}\frac{M^{208}}{A^{156}}
 =\frac{A^{306}}{M^{485}}=z_k.
\]
The third vector $(5,7)$ is the sum of the first two, and its tangent
$(615,-283)$ is their sum.  The corresponding constant is
$C_{(5,7)}=C_\ell C_r$, illustrating the rank-two syzygy directly.

\begin{statusbox}{\exactcheck.  Every displayed identity is exact integer algebra.  The
example is included to make the constant-covolume geometry concrete, not as a statistical
experiment.}
\end{statusbox}

\section{Consequences for the surviving mask architecture}
\label{sec:consequences}

\subsection{What is now closed}

The following branches of the shared structural-face program can be retired.

\begin{enumerate}
\item \textbf{Fixed or sublinear total degree.}
Corollary~\ref{cor:sublinear-no-face} shows that no stable denominator prime has a tied
geometric face once $d_k+e_k=o(q_k)$.  Any such construction is coefficient-driven from the
start.

\item \textbf{A single prime supplying geometric cancellation.}
Theorems~\ref{thm:edge-evaluation} and~\ref{thm:one-prime-normal-form} show that a tied edge
either contributes only the bounded rational value $P(C_\ell^{\pm1})$, or contains the exact
tangent binomial and can be divided out.  Repeating the division removes every one-prime
geometric tie.

\item \textbf{Area-only Newton polygon arguments.}
Theorem~\ref{thm:tangent-determinant} produces arbitrarily long pairs of tangent edges with
constant determinant.  A Newton polygon may be extremely thin.  Any lower bound based only
on area, mixed volume, or the number of vertices misses the critical configurations.

\item \textbf{Prime-count independence.}
The Pl\"ucker relations show that all tangent constants have multiplicative rank at most two.
More support primes do not automatically yield more independent rational evaluations.

\item \textbf{A new irrationality measure from local depth ratios.}
Equation~\eqref{eq:rho} is a fractional-linear transform of $\thetaq$, and
\eqref{eq:transported-error} transports the same convergents.  Baker-type bounds for the
local logarithm ratios are equivalent in strength to bounds for the original two logarithms.
\end{enumerate}

These are negative results about a method, not about the truth of the conjecture.  They make
the remaining method sharply identifiable.

\subsection{The first genuinely unresolved mechanism}

After one-prime reduction, a mask can still obtain deep denominator compression in three
ways:

\begin{enumerate}
\item its normal-form coefficients carry large structural-prime divisibility;
\item units from different normal-form monomials satisfy deep, output-specific congruences;
\item reductions demanded by two or more nonproportional prime slopes interact so that a
coefficient expensive in one local basis is cheap in another.
\end{enumerate}

The third possibility is the new focal point.  The tangent lattice has fixed index, so
simultaneous reduction has only finitely many exponent residue classes.  But the basis is
ill-conditioned at order $q_k$, so the rational coefficient multipliers are exactly large
enough to absorb the apparent denominator saving.  A proof or definitive no-go theorem must
control this coefficient transformation place by place.

\subsection{Relation to the Matrix Coefficient frontier}

Dasgupta and Kakde's Matrix Coefficient Conjecture framework concerns sparse auxiliary
polynomials and ranks of logarithm matrices \cite{DasguptaKakde2026}.  The tangent binomials
here are exceptionally sparse---two monomials---but their exponents are long and depend on the
hypothetical output valuations.  They therefore provide exact sparse relations at one
near-unit pair without producing the global fewnomial required by the matrix-coefficient
argument.

The new normal form can nevertheless serve as preprocessing for that program.  Any proposed
fewnomial supported on the near-unit orbit should first be reduced modulo the tangent
binomials.  Monomials in the same tangent coset are redundant at the evaluation point; after
reduction, the true sparse complexity is the number of residue classes together with the
height of their transformed coefficients.  This separates support count from coefficient
height in an exact way.

\subsection{Relation to weight-two transcendence}

The fundamental counterexample relation
\[
 \log M\log3-\log A\log2=0
\]
has weight two in logarithms.  The local constants do not lower that weight: by
\eqref{eq:logC}, their logarithms are rational-linear expressions in $\log M,\log A$, while
ratios of those logarithms are fractional-linear transforms of $\thetaq$.  Baker--W\"ustholz
and the Laurent--Mignotte--Nesterenko two-logarithm bounds remain weight-one tools
\cite{BakerWustholz1993,LaurentMignotteNesterenko}.  The cross-prime normal form may organize
the coefficients of a future weight-two argument, but it does not itself create the missing
transcendence theorem.

\section{A sharpened research program}
\label{sec:program}

\subsection{Priority A: cross-prime normal-form contraction}

\begin{target}[Cross-prime normal-form contraction]
\label{target:cross-prime-contraction}
Fix two nonproportional stable denominator slopes $\ell,r$.  Develop a simultaneous integral
normal form modulo the two tangent binomials and prove one of the following mutually decisive
statements.
\begin{enumerate}
\item \emph{Negative closure:} every nonzero mask of contact $R$ and subcritical original
coefficient height retains reduced denominator
\[
 \log\den F(z_k,z_{k+1})
 \ge c_{\ell r}\,q_k\,\operatorname{diam}(\supp F)
   -O(R\log q_k+\log\lVert F\rVert_1)
\]
for some $c_{\ell r}>0$; or
\item \emph{Constructive breakthrough:} exhibit masks for which simultaneous reduction
captures a $1-o(1)$ fraction of the denominator depth at every stable denominator prime while
keeping transformed coefficient height $O(R\log q_k)$.
\end{enumerate}
\end{target}

The first outcome would eliminate the surviving mask architecture.  The second, together
with the established archimedean contact estimates in the unified report, would be a credible
route to the conjecture.  The present results show that no intermediate one-prime geometric
phenomenon can decide the issue.

\subsection{Priority B: exact coefficient modules on thin tangent cells}

For a pair with $|\Delta_{\ell r}|=D$, choose a fundamental domain
$\mathcal R\subset\Z^2$ of $D$ exponent residues.  Every Laurent polynomial has a unique
representation at the evaluation point
\begin{equation}
 F(z_k,z_{k+1})
 =\sum_{\rho\in\mathcal R}Q_{\rho,k}(C_\ell,C_r)
   z_k^{\rho_1}z_{k+1}^{\rho_2},
 \label{eq:simultaneous-normal}
\end{equation}
where the $Q_{\rho,k}$ are Laurent polynomials with rational coefficients obtained by lattice
reduction.  The exact computational problem is to determine the Smith profile of the map
from original integer coefficients to the numerator-denominator pairs of all
$Q_{\rho,k}$.

The long tangent basis makes generic Hermite or nearest-plane reduction misleading.  The
right basis is supplied by the exact short vectors \eqref{eq:short-e1}--\eqref{eq:short-e2}.
A practical algorithm should:
\begin{enumerate}
\item reduce exponents by those short vectors rather than by a floating-point lattice basis;
\item keep numerator and denominator prime factorizations separately;
\item quotient common content after every move;
\item return an exact sparse matrix whose Smith form records unavoidable coefficient taxes.
\end{enumerate}

This is finite integer linear algebra for every fixed output valuation configuration.

\subsection{Priority C: a two-prime initial-form incompatibility theorem}

At one prime, the vanishing initial form is the rational-root condition
$P(C_\ell^{\pm1})=0$.  At two primes, different edges give polynomial conditions at
$C_\ell$ and $C_r$.  Because the constants are multiplicatively independent when the slopes
are nonproportional, a promising target is:

\begin{target}[Primitive simultaneous edge cancellation]
\label{target:edge-incompatibility}
Show that if a primitive integer polynomial mask has deeply vanishing initial forms at two
nonproportional structural primes, then either
\begin{enumerate}
\item it lies in the ideal generated by the two tangent binomials, hence vanishes exactly at
$(z_k,z_{k+1})$; or
\item the content or height of its transformed coefficient vector is at least
$\exp(cq_kD)$, where $D$ is the relevant support diameter.
\end{enumerate}
\end{target}

A proof would likely combine rational-root divisibility, resultants of the two edge
polynomials, and the constant-index Smith form.  Theorem~\ref{thm:resultant-identities}
shows that the basic binomial resultant alone is exactly critical; an improvement must use
primitivity of the full coefficient system or a third place.

\subsection{Priority D: exploit a third place without assuming a third rank}

The Pl\"ucker relation prevents three tangent constants from being multiplicatively
independent.  A third prime can still help through local incompatibility: a coefficient
transformation that is integral at $\ell$ and $r$ may have a forced denominator at $s$.
This is an adelic Smith problem rather than a rank problem.

A concrete finite target is to choose three slope classes and compute the intersection of the
three integral coefficient lattices obtained from their one-prime normal forms.  If the
intersection has covolume growing like $\exp(cq_kD)$ after primitive content is removed, then
the product formula supplies the missing global tax.  If its covolume stays bounded, the
computation should expose an explicit three-prime auxiliary mask.

\subsection{Priority E: classify the one-denominator-slope case}

It may happen that all primes in $\mathcal P_-$ share one primitive valuation slope.  Then a
single tangent binomial removes geometric ties at every denominator prime simultaneously.
The mask problem becomes a coefficient-congruence problem in one normal form, without the
cross-prime conditioning barrier on the denominator side.

This special case deserves separate treatment.  Two possible outcomes are:
\begin{enumerate}
\item prove a primitive coefficient-descent theorem and eliminate the case; or
\item show that the numerator support must then contain at least two nonproportional slopes
and transfer the argument through a reciprocal or Laurent mask.
\end{enumerate}
The slope partition is finite and computable from the output valuation matrix, so the case
split is exact.

\subsection{Priority F: controlled computational reconnaissance}

The computational search should no longer enumerate arbitrary high-contact masks.  The new
geometry suggests the following restricted families:
\begin{enumerate}
\item unimodular tangent triangles when $|\Delta_{\ell r}|=1$;
\item fundamental parallelograms for $|\Delta_{\ell r}|>1$;
\item one or two Minkowski sums of those cells to allow higher contact at $(1,1)$;
\item coefficient vectors primitive after every one-prime tangent division;
\item simultaneous local constraints expressed by exact Smith normal forms, not floating
point valuations.
\end{enumerate}

For each family, the output should be either a candidate mask or a Farkas/Smith certificate
showing that the required local divisibility forces coefficient height above the
archimedean budget.  Searches on control pairs must be run in parallel: any certificate that
also excludes the integer controls is diagnosing a method, not the conjecture.

\section{Lean-oriented formalization blueprint}
\label{sec:lean}

The new core is elementary integer algebra and is well suited to formalization.  The suggested
module order is intentionally dependency-light.

\subsection{Module 1: adjacent convergent matrices}

Define an integer matrix
\[
 B(p,q,p',q')=\begin{pmatrix}-p&-p'\\q&q'\end{pmatrix}
\]
and assume $p'q-pq'=\eps$ with $\eps^2=1$.  Prove:
\begin{enumerate}
\item determinant $\eps$;
\item explicit integral inverse;
\item preservation of the ideal generated by a row vector's coordinates;
\item preservation of all $2\times2$ minors under right multiplication.
\end{enumerate}
Suggested theorem names:
\begin{verbatim}
adjacentMatrix_det
adjacentMatrix_mem_GL2Z
spanIdeal_mul_adjacentMatrix
row_gcd_mul_adjacentMatrix
minor_mul_adjacentMatrix
\end{verbatim}

\subsection{Module 2: local depth transport}

With $u_k=qn-pm$ and $u_{k+1}=q'n-p'm$, prove:
\begin{verbatim}
localDepthPair_eq_mul
localDepth_gcd
localDepth_crossDet
localDepth_q_transport
localDepth_p_transport
\end{verbatim}
The first three should be proved from the matrix module, avoiding repeated ring expansions.

\subsection{Module 3: primitive tangent and cross-ratio}

For positive $s=-u_k$, $s'=-u_{k+1}$ and $h=\gcd(m,n)$, define
$a=s'/h$, $b=s/h$.  Prove:
\begin{verbatim}
primitiveTangent_coprime
equalPoleWeight_iff
nearUnit_tangent_crossRatio
crossRatio_primeUnit
crossRatio_gt_one
\end{verbatim}
The cross-ratio equality can be formulated first in a commutative group with integer powers,
then specialized to positive rationals.

\subsection{Module 4: edge polynomial and normal form}

Formalize the map
\[
 U^iV^j\longmapsto
 N^{\lfloor i/a\rfloor}D^{T-\lfloor i/a\rfloor}
 U^{i\bmod a}V^{j+b\lfloor i/a\rfloor}.
\]
Prove evaluation, degree, support-injectivity, and coefficient $\ell^1$ bounds.  Suggested
statements:
\begin{verbatim}
tangentNormalForm_eval
tangentNormalForm_natDegree_U_lt
tangentNormalForm_natDegree_V_le
tangentNormalForm_l1Norm_le
tangentNormalForm_poleWeight_injective
tangentNormalForm_eq_zero_iff_dvd
\end{verbatim}
The divisibility theorem may be easiest over $\Q[V][U]$, followed by a primitive-polynomial
Gauss lemma.

\subsection{Module 5: cross-prime tangent lattice}

Formalize:
\begin{verbatim}
tangent_crossDet
tangentLattice_index
tangent_shortVector_e1
tangent_shortVector_e2
nearUnit_resultant_first
nearUnit_resultant_second
plucker_crossRatio
\end{verbatim}
No analysis or transcendence theory is needed.  The asymptotic length statements can remain
paper-level initially; the exact algebra is the reusable core.

\subsection{Formalization payoff}

A kernel-checked implementation would provide three durable guarantees:
\begin{enumerate}
\item every future mask search can quotient exact one-prime redundancies before optimization;
\item all sign conventions in cross-prime resultants become mechanically stable;
\item computed Smith certificates can be replayed against a small, trusted algebraic API.
\end{enumerate}
This is a higher-value formalization target than extending the finite interval scan by another
doubling, because it changes the search space rather than only the verified range.

\section{Exact verification script}
\label{sec:script}

The following Python program uses exact integer and rational arithmetic for every asserted
identity.  The continued-fraction partial quotients are supplied as data; floating-point
arithmetic is not used in the checks.  The program verifies:
\begin{enumerate}
\item adjacency and determinant signs for the sample convergents;
\item exact gcd preservation for local depths;
\item the cross-ratio exponents in $A$ and $M$;
\item the cross-prime tangent determinant and short-vector identities;
\item evaluation preservation for the canonical one-prime normal form on an independent
synthetic rational point.
\end{enumerate}

\begin{lstlisting}[style=pythonstyle,caption={Exact checks for transported local depths and tangent normal forms.}]
from __future__ import annotations

from dataclasses import dataclass
from fractions import Fraction
from math import gcd
from typing import Iterable


@dataclass(frozen=True)
class ConvergentPair:
    p: int
    q: int
    p1: int
    q1: int

    @property
    def eps(self) -> int:
        value = self.p1 * self.q - self.p * self.q1
        if value not in (-1, 1):
            raise ValueError(f"not adjacent: determinant is {value}")
        return value


@dataclass(frozen=True)
class LocalData:
    m: int
    n: int
    h: int
    s: int
    s1: int
    a: int
    b: int


def continued_fraction_convergents(partials: Iterable[int]) -> list[tuple[int, int]]:
    p_m2, p_m1 = 0, 1
    q_m2, q_m1 = 1, 0
    result: list[tuple[int, int]] = []
    for digit in partials:
        p = digit * p_m1 + p_m2
        q = digit * q_m1 + q_m2
        result.append((p, q))
        p_m2, p_m1 = p_m1, p
        q_m2, q_m1 = q_m1, q
    return result


def local_data(pair: ConvergentPair, m: int, n: int) -> LocalData:
    h = gcd(m, n)
    s = m * pair.p - n * pair.q
    s1 = m * pair.p1 - n * pair.q1
    if s <= 0 or s1 <= 0:
        raise ValueError("the chosen vector is not a stable denominator vector here")
    assert gcd(s, s1) == h
    return LocalData(m=m, n=n, h=h, s=s, s1=s1, a=s1 // h, b=s // h)


def det(x: tuple[int, int], y: tuple[int, int]) -> int:
    return x[0] * y[1] - x[1] * y[0]


def verify_pair(pair: ConvergentPair, left: LocalData, right: LocalData) -> None:
    tl = (left.a, -left.b)
    tr = (right.a, -right.b)
    normalized_output_det = (
        left.m * right.n - left.n * right.m
    ) // (left.h * right.h)
    assert det(tl, tr) == pair.eps * normalized_output_det

    signed_delta = det(tl, tr)
    e1 = (-right.b * tl[0] + left.b * tr[0],
          -right.b * tl[1] + left.b * tr[1])
    e2 = (-right.a * tl[0] + left.a * tr[0],
          -right.a * tl[1] + left.a * tr[1])
    assert e1 == (signed_delta, 0)
    assert e2 == (0, signed_delta)


def cross_ratio_exponents(pair: ConvergentPair, data: LocalData) -> tuple[int, int, int, int]:
    # z_k^a / z_{k+1}^b = A^Aexp M^Mexp.
    Aexp = pair.q * data.a - pair.q1 * data.b
    Mexp = -pair.p * data.a + pair.p1 * data.b
    assert Aexp == pair.eps * (data.m // data.h)
    assert Mexp == -pair.eps * (data.n // data.h)
    return Aexp, Mexp, data.a, data.b


def normal_form_terms(
    terms: dict[tuple[int, int], int], a: int, b: int, N: int, D: int
) -> dict[tuple[int, int], int]:
    if gcd(N, D) != 1 or min(a, b, N, D) <= 0:
        raise ValueError("positive primitive data required")
    degree_u = max((i for (i, _j) in terms), default=0)
    T = degree_u // a
    out: dict[tuple[int, int], int] = {}
    for (i, j), coefficient in terms.items():
        t, r = divmod(i, a)
        key = (r, j + b * t)
        out[key] = out.get(key, 0) + coefficient * (N**t) * (D ** (T - t))
    return {key: value for key, value in out.items() if value != 0}


def eval_poly(terms: dict[tuple[int, int], int], u: Fraction, v: Fraction) -> Fraction:
    return sum((Fraction(c) * (u**i) * (v**j) for (i, j), c in terms.items()), Fraction(0))


def verify_normal_form() -> None:
    # An exact synthetic point with D U^a = N V^b.
    a, b, N, D = 3, 2, 5, 7
    # Here U=V=N/D works because a=b+1.
    u = Fraction(5, 7)
    v = Fraction(5, 7)
    assert D * (u**a) == N * (v**b)

    terms = {
        (0, 0): 11,
        (1, 2): -3,
        (3, 0): 7,
        (4, 1): 2,
        (7, 0): -5,
    }
    transformed = normal_form_terms(terms, a, b, N, D)
    T = max(i for i, _j in terms) // a
    assert eval_poly(transformed, u, v) == (D**T) * eval_poly(terms, u, v)
    assert all(i < a for i, _j in transformed)


def main() -> None:
    # Initial partial quotients of log_2(3).
    partials = [1, 1, 1, 2, 2, 3, 1, 5, 2, 23, 2, 2, 1, 1, 55]
    conv = continued_fraction_convergents(partials)
    (p, q), (p1, q1) = conv[7], conv[8]
    pair = ConvergentPair(p=p, q=q, p1=p1, q1=q1)
    assert (p, q, p1, q1, pair.eps) == (485, 306, 1054, 665, -1)

    d1 = local_data(pair, 2, 3)
    d2 = local_data(pair, 3, 4)
    d3 = local_data(pair, 5, 7)
    assert (d1.s, d1.s1, d1.a, d1.b) == (52, 113, 113, 52)
    assert (d2.s, d2.s1, d2.a, d2.b) == (231, 502, 502, 231)
    assert (d3.a, -d3.b) == (d1.a + d2.a, -(d1.b + d2.b))

    for data in (d1, d2, d3):
        cross_ratio_exponents(pair, data)
    verify_pair(pair, d1, d2)
    verify_normal_form()

    print("all exact transport, cross-ratio, determinant, and normal-form checks passed")


if __name__ == "__main__":
    main()
\end{lstlisting}

The expected output is:
\begin{verbatim}
all exact transport, cross-ratio, determinant, and normal-form checks passed
\end{verbatim}

\begin{statusbox}{\exactcheck.  The script was executed successfully while preparing this
report.  It is a consistency check, not a substitute for the proofs.}
\end{statusbox}

\section{Status matrix and dependency map}
\label{sec:status}

\begin{longtable}{@{}p{0.29\textwidth}p{0.17\textwidth}p{0.46\textwidth}@{}}
\toprule
Statement & Status & Dependencies and role\\
\midrule
\endfirsthead
\toprule
Statement & Status & Dependencies and role\\
\midrule
\endhead
Exact unimodular transport, Theorem~\ref{thm:unimodular-transport}
& \provedhere
& Continued-fraction determinant only; foundational.\\

Consecutive gcd invariance, Corollary~\ref{cor:gcd-invariance}
& \provedhere
& Unimodularity; supplies primitive tangents.\\

Full Smith and minor invariance, Corollary~\ref{cor:smith-invariance}
& \provedhere
& Right $\GL_2(\Z)$ equivalence; local-rank audit.\\

Cross-prime determinant, Theorem~\ref{thm:cross-determinant}
& \provedhere
& Determinant multiplicativity; fixed slope classes.\\

Full tangent-matrix equivalence, Theorem~\ref{thm:tangent-matrix-equivalence}
& \provedhere
& A common right $\GL_2(\Z)$ transform preserves all tangent Smith data.\\

Stable sign split and linear depths
& \provedhere
& Convergent error and $\log A-\thetaq\log M=0$.\\

Fixed tangent cross-ratio, Theorem~\ref{thm:fixed-crossratio}
& \provedhere
& Transport identities; central local invariant.\\

Projective transport of local convergents
& \provedhere
& Fractional-linear algebra; no new irrationality measure.\\

Primitive Newton tangent, Theorem~\ref{thm:primitive-tangent}
& \provedhere
& Exact gcd; classifies all tied monomials.\\

Linear face-width and sublinear no-face results
& \provedhere
& Primitive tangent length; closes a search branch.\\

Multi-prime perimeter tax, Theorem~\ref{thm:perimeter-tax}
& \provedhere
& Convex Newton boundary; additive width cost.\\

Edge-polynomial evaluation, Theorem~\ref{thm:edge-evaluation}
& \provedhere
& Fixed cross-ratio; reduces a face to $P(C_\ell)$.\\

Rational edge valuation bound
& \provedhere
& Elementary integer numerator bound.\\

Canonical one-prime normal form, Theorem~\ref{thm:one-prime-normal-form}
& \provedhere
& Polynomial division by the tangent binomial.\\

Constant-index tangent lattice, Corollary~\ref{cor:constant-index}
& \provedhere
& Cross-prime determinant; explains thin cells.\\

Exact two-prime resultants, Theorem~\ref{thm:resultant-identities}
& \provedhere
& Short-vector identities; exposes conditioning.\\

Three-prime Pl\"ucker syzygy, Theorem~\ref{thm:plucker}
& \provedhere
& Rank-two exponent lattice; many-prime ceiling.\\

Synthetic calibration and Python script
& \exactcheck
& Independent sign and normal-form checks.\\

Gelfond--Schneider transcendence of a counterexample
& \knownstatus
& Classical background; not needed for transport algebra.\\

Four-exponentials implication
& \knownstatus
& Explains the global barrier; conjectural input.\\

Cross-prime normal-form contraction
& \openstatus
& Principal sharpened target left by this report.\\

Primitive simultaneous edge cancellation theorem
& \openstatus
& Would convert constant-index geometry into a height tax.\\
\bottomrule
\end{longtable}

The logical dependency graph of the new exact core is short:
\[
\boxed{\text{adjacent convergents}}
\Longrightarrow
\boxed{\text{unimodular depth transport}},
\]
then
\[
\boxed{\text{gcd/Smith invariance}}
\qquad
\boxed{\text{fixed cross-ratios}}
\qquad
\boxed{\text{cross-prime determinants}},
\]
followed by
\[
\boxed{\text{fixed cross-ratio}}
\Longrightarrow
\boxed{\text{edge polynomial}}
\Longrightarrow
\boxed{\text{one-prime normal form}},
\]
and
\[
\boxed{\text{cross-prime determinant}}
\Longrightarrow
\boxed{\text{constant-index tangent lattice}}
\Longrightarrow
\boxed{\text{critical conditioning barrier}}.
\]
No transcendence estimate enters this new core.

\section{Final assessment}
\label{sec:conclusion}

The \AEname{} conjecture remains open.  The work reported here does not cross the
four-exponentials barrier, and it would be misleading to present the new identities as a
near-complete proof.  They are valuable for a different reason: they settle the internal
geometry of the most concrete surviving mask target in the supplied report.

The decisive facts are exact.

\begin{enumerate}
\item Adjacent local depths are not an irregular moving pair.  They are a unimodular image of
the fixed output-valuation vector, so gcds, minors, Smith data, and slope classes are rigid.
\item The primitive direction of every tied structural-prime face is long, but its monomial
ratio is a fixed rational number.  A face is therefore a one-variable polynomial at a fixed
rational point.
\item Nonzero face values have a direct height bound.  Zero face values are tangent-binomial
factors.  One-prime geometric cancellation can be divided out canonically at low relative
cost.
\item Different prime slopes produce long, almost parallel tangent vectors of constant
determinant.  Simultaneous quotienting has finite exponent index but critical coefficient
conditioning.  Exact resultants reproduce the original near-unit height, so no contradiction
comes from total height alone.
\item Three or more prime constants remain in a rank-two multiplicative group and satisfy
explicit Pl\"ucker syzygies.  The remaining gain, if it exists, must be adelic and
coefficient-sensitive rather than rank-theoretic.
\end{enumerate}

Accordingly, the frontier has moved from the broad question ``can shared Newton faces cancel
deeply?'' to a narrower one:

\begin{center}
\begin{minipage}{0.88\textwidth}
\emph{Can a primitive integer coefficient system be simultaneously cheap in two or more
nonproportional, constant-index but $q$-ill-conditioned tangent normal forms?}
\end{minipage}
\end{center}

A negative answer with an effective height lower bound closes the mask architecture.  A
positive answer produces exactly the kind of output-specific denominator compression that
the larger determinant program has been missing.  Either outcome is mathematically useful,
and both are now finite enough to support exact Smith-form computation and Lean
formalization.

\appendix

\section{Sign audit for the cross-ratio and resultant formulas}
\label{app:signs}

Because the convergent determinant alternates, sign errors are easy to introduce.  This
appendix records every convention in one place.

Let
\[
 \eps_k=p_{k+1}q_k-p_kq_{k+1}\in\{\pm1\},
 \qquad
 u_k=q_kn-p_km,
 \qquad
 s_k=-u_k=p_km-q_kn.
\]
Then
\begin{align*}
 q_{k+1}u_k-q_ku_{k+1}&=\eps_km,\\
 p_{k+1}u_k-p_ku_{k+1}&=\eps_kn,\\
 q_ks_{k+1}-q_{k+1}s_k&=\eps_km,\\
 p_{k+1}s_k-p_ks_{k+1}&=-\eps_kn.
\end{align*}
With $h=\gcd(m,n)$, $a=s_{k+1}/h$, $b=s_k/h$,
\begin{align*}
 \frac{z_k^a}{z_{k+1}^b}
 &=A^{(q_ks_{k+1}-q_{k+1}s_k)/h}
   M^{(p_{k+1}s_k-p_ks_{k+1})/h}\\
 &=A^{\eps_km/h}M^{-\eps_kn/h}
 =\left(A^{m/h}M^{-n/h}\right)^{\eps_k}.
\end{align*}

For two primes, put $t_\ell=(a_\ell,-b_\ell)$ and
$t_r=(a_r,-b_r)$.  Then
\[
 \Delta=\det(t_\ell,t_r)=-a_\ell b_r+a_rb_\ell.
\]
Consequently
\[
 -b_rt_\ell+b_\ell t_r=(\Delta,0),
 \qquad
 -a_rt_\ell+a_\ell t_r=(0,\Delta),
\]
which gives
\[
 z_k^\Delta=(C_\ell^{-b_r}C_r^{b_\ell})^{\eps_k},
 \qquad
 z_{k+1}^\Delta=(C_\ell^{-a_r}C_r^{a_\ell})^{\eps_k}.
\]
These formulas remain correct for negative $\Delta$ because rational powers here are integer
powers and inversion is allowed.

\section{Full proof of the simultaneous slope determinant formula}
\label{app:detproof}

For completeness, expand Theorem~\ref{thm:tangent-determinant} directly.  Write
$s_{\ell,k}=p_km_\ell-q_kn_\ell$ and similarly for $r$.  Then
\begin{align*}
&s_{\ell,k}s_{r,k+1}-s_{\ell,k+1}s_{r,k}\\
&=(p_km_\ell-q_kn_\ell)(p_{k+1}m_r-q_{k+1}n_r)
 -(p_{k+1}m_\ell-q_{k+1}n_\ell)(p_km_r-q_kn_r).
\end{align*}
The $m_\ell m_r$ and $n_\ell n_r$ terms cancel.  The coefficient of $m_\ell n_r$ is
\[
 -p_kq_{k+1}+p_{k+1}q_k=\eps_k,
\]
while the coefficient of $n_\ell m_r$ is $-\eps_k$.  Hence
\[
 s_{\ell,k}s_{r,k+1}-s_{\ell,k+1}s_{r,k}
 =\eps_k(m_\ell n_r-n_\ell m_r).
\]
Divide by $h_\ell h_r$ and use
$t_{\ell,k}=(s_{\ell,k+1}/h_\ell,-s_{\ell,k}/h_\ell)$.

The matrix proof in the main text is shorter and scales to all minors.  The expansion is
included as an independent sign audit.

\section{Normal form as a quotient-ring statement}
\label{app:quotient}

Let $K$ be a field, $c\in K^\times$, and $a,b$ positive.  In
$K[U,V]/(U^a-cV^b)$ every class has a unique representative of $U$-degree $<a$ because the
relation is monic in $U$.  The representative is obtained by repeated replacement
$U^a\mapsto cV^b$.  Taking $K=\Q$ and $c=N/D$ gives the rational normal form; multiplying by
$D^T$ gives the integral representative of Theorem~\ref{thm:one-prime-normal-form}.

If $\gcd(a,b)=1$, the binomial $U^a-cV^b$ is irreducible over an algebraic closure after a
nonzero scalar extension and defines a prime Laurent ideal.  Irreducibility is not needed for
uniqueness of the $U$-remainder, but it clarifies the geometry: the tangent relation defines a
one-dimensional torus coset, and the normal form is ordinary coordinate reduction on that
coset.

For two nonproportional tangent relations, the exponent lattice has rank two and finite index
$D=|\Delta|$.  The corresponding Laurent quotient is finite-dimensional of dimension $D$
over a suitable coefficient field, with one basis monomial for each exponent residue class.
The hard part is not dimension but integral height: expressing a fixed residue representative
uses a badly conditioned long basis, as quantified in Section~\ref{sec:crossprime}.

\section{A compact checklist for future auxiliary masks}
\label{app:checklist}

Every proposed bivariate near-unit mask should pass the following audit before large-scale
computation.

\begin{enumerate}
\item \textbf{Degree scale.}  Is $d+e$ at least the primitive tangent length for every prime
where a tied face is claimed?  If not, the face cannot exist.
\item \textbf{Slope deduplication.}  Group support primes by primitive valuation vector.
Several primes in one class impose the same tangent direction and multiplicatively dependent
cross-ratios.
\item \textbf{Edge reduction.}  Replace each tied face by its one-variable polynomial
$P(C_\ell^{\pm1})$.  If nonzero, charge the bound
\eqref{eq:edge-extra-bound}; if zero, divide the tangent binomial.
\item \textbf{Normal-form primitivity.}  Remove common coefficient content after every
one-prime division.  Apparent local gain carried entirely by coefficient content is descent,
not cancellation.
\item \textbf{Cross-prime index.}  Compute $\Delta_{\ell r}$.  Use a fundamental domain of
the tangent lattice rather than an arbitrary rectangular monomial box.
\item \textbf{Conditioning tax.}  Track powers of $C_\ell,C_r$ exactly.  Do not estimate them
independently before reducing the rational product; Theorem~\ref{thm:resultant-identities}
shows large placewise cancellations.
\item \textbf{Third-place audit.}  After simultaneous reduction at two primes, factor every
transformed coefficient and inspect the other support primes.  A useful gain must survive the
product formula.
\item \textbf{Control test.}  Run the same construction on integer-control pairs.  A method
that closes controls as well as candidates has proved a no-go theorem about itself.
\item \textbf{Exact certificate.}  Return a Smith/Farkas certificate or a literal auxiliary
polynomial.  Floating-point evidence is not enough at the critical scale.
\end{enumerate}

\begin{thebibliography}{99}

\bibitem{AE1944}
L.~Alaoglu and P.~Erd\H{o}s,
\newblock \emph{On highly composite and similar numbers},
\newblock Transactions of the American Mathematical Society \textbf{56} (1944), 448--469.
\newblock DOI: \href{https://doi.org/10.1090/S0002-9947-1944-0011087-2}{10.1090/S0002-9947-1944-0011087-2}.

\bibitem{Baker1975}
A.~Baker,
\newblock \emph{Transcendental Number Theory},
\newblock Cambridge University Press, 1975.

\bibitem{BakerWustholz1993}
A.~Baker and G.~W\"ustholz,
\newblock \emph{Logarithmic forms and group varieties},
\newblock Journal f\"ur die reine und angewandte Mathematik \textbf{442} (1993), 19--62.
\newblock DOI: \href{https://doi.org/10.1515/crll.1993.442.19}{10.1515/crll.1993.442.19}.

\bibitem{DasguptaKakde2026}
S.~Dasgupta and M.~Kakde,
\newblock \emph{Ranks of matrices of logarithms of algebraic numbers II: The Matrix
Coefficient Conjecture},
\newblock Advances in Mathematics \textbf{487} (2026), Article 110753.
\newblock DOI: \href{https://doi.org/10.1016/j.aim.2025.110753}{10.1016/j.aim.2025.110753};
arXiv:2408.08178.

\bibitem{LangTranscendental}
S.~Lang,
\newblock \emph{Introduction to Transcendental Numbers},
\newblock Addison--Wesley, 1966.

\bibitem{LaurentMignotteNesterenko}
M.~Laurent, M.~Mignotte, and Yu.~Nesterenko,
\newblock \emph{Formes lin\'eaires en deux logarithmes et d\'eterminants d'interpolation},
\newblock Journal of Number Theory \textbf{55} (1995), 285--321.

\bibitem{Unified6}
V.~Reshetnikov,
\newblock \emph{The Alaoglu--Erd\H{o}s Integer-Exponent Conjecture: Unified Report},
\newblock working report, version 6, August 2026.

\bibitem{Waldschmidt2023}
M.~Waldschmidt,
\newblock \emph{The Four Exponentials Problem and Schanuel's Conjecture},
\newblock in \emph{Mathematics Going Forward}, Lecture Notes in Mathematics, 2023,
pp.~579--592.
\newblock DOI: \href{https://doi.org/10.1007/978-3-031-12244-6_39}{10.1007/978-3-031-12244-6\_39}.

\end{thebibliography}

\end{document}
